\documentclass[11pt,a4paper]{article}

\usepackage[T1]{fontenc}
\usepackage[utf8]{inputenc}
\usepackage[english]{babel}
\usepackage[a4paper,margin=2.4cm]{geometry}
\usepackage{amsmath,amssymb,amsthm}
\usepackage{booktabs}
\usepackage{longtable}
\usepackage{array}
\usepackage{enumitem}
\usepackage{microtype}
\usepackage[hidelinks]{hyperref}

\newcolumntype{L}[1]{>{\raggedright\arraybackslash}p{#1}}
\newcolumntype{C}[1]{>{\centering\arraybackslash}p{#1}}
\newcommand{\mpart}[1]{\smallskip\noindent\textbf{#1}\ \ignorespaces}
\setlist[itemize]{noitemsep,topsep=3pt,leftmargin=*}
\setlist[enumerate]{noitemsep,topsep=3pt,leftmargin=*}
\setlength{\parskip}{2pt}
\renewcommand{\arraystretch}{1.15}

\theoremstyle{definition}
\newtheorem{definition}{Definition}
\newtheorem{example}[definition]{Example}
\theoremstyle{plain}
\newtheorem{theorem}{Theorem}
\newtheorem{lemma}[theorem]{Lemma}
\newtheorem{proposition}[theorem]{Proposition}
\newtheorem{corollary}[theorem]{Corollary}
\theoremstyle{remark}
\newtheorem*{remark}{Remark}

\newcommand{\R}{\mathbb{R}}
\newcommand{\ip}[2]{\langle #1,\, #2\rangle}
\newcommand{\norm}[1]{\lVert #1\rVert}
\newcommand{\proj}{P}
\DeclareMathOperator{\conv}{conv}
\DeclareMathOperator{\aff}{aff}
\DeclareMathOperator{\cone}{cone}
\DeclareMathOperator{\ri}{ri}
\DeclareMathOperator{\cl}{cl}
\DeclareMathOperator{\bd}{bd}
\DeclareMathOperator{\dist}{d}
\DeclareMathOperator{\prox}{prox}
\DeclareMathOperator{\spn}{span}
\DeclareMathOperator*{\argmin}{arg\,min}
\DeclareMathOperator*{\argmax}{arg\,max}
\newcommand{\Simplex}{\Delta}
\newcommand{\Nc}{N}
\newcommand{\one}{\mathbf{1}}

\title{\textbf{Module M1, Session by Session:\\ Convex Sets, Separation and Farkas' Lemma}\\[3pt]
\large MMC007 --- Optimization}
\author{MSc in Mathematical and Computational Modeling\\
IMCA --- Faculty of Sciences, Universidad Nacional de Ingenier\'ia}
\date{Semester 2026--2}

\begin{document}

\maketitle

\section{Purpose and place in the plan}

M1 is the first module of the MMC007 reading plan and the only one with no prerequisites. It
covers the first item of the syllabus --- ``properties of convex sets'' --- and it answers a
question that item normally leaves unasked: \emph{why does one geometric fact about convex sets
generate every duality theorem in the course?}

The answer is that a point outside a closed convex set can be strictly separated from it by a
hyperplane; that the separating vector is not an abstract existence but a vector a projection
computes; and that every certificate of optimality or of infeasibility in the rest of the plan
is an instance of that one construction. Three questions become askable at once, and this
module answers all three:

\begin{enumerate}
\item Given a point and a convex set, how is the separating hyperplane \emph{produced}?
Answered by the projection theorem (S3) and its two-line corollary (S4).
\item When does separation \emph{fail}? Answered by S2 and by the capstone of S7: it fails
without closedness, and every duality gap in the course is a failure of exactly this kind.
\item What does separation give when the set is a finitely generated cone? Answered by Farkas'
lemma (S6), which is the single result the rest of the plan quotes most often.
\end{enumerate}

Every later module inherits something specific from M1:

\begin{itemize}
\item The \emph{projection theorem} of S3 reappears in M8 as the stationarity condition
$x^\star=\proj_C(x^\star-t\nabla f(x^\star))$ of the projected gradient method, and its firm
nonexpansiveness is what makes the proximal iterates of M9 Fej\'er monotone. It is also, with
$C$ an affine subspace, the projection theorem MMC201's module M2 uses for least squares: one
theorem, two courses, two guises.
\item The \emph{separation theorem} of S4 is the engine of M4. The proof of strong duality
under Slater's condition separates the point $(0,p^\star)$ from the perturbation epigraph, and
it is the same proof written here.
\item \emph{Farkas' lemma}, proved in S6, is repeated in M4 with the tangent cone in place of
$\cone(A)$ --- and the result is the Karush--Kuhn--Tucker theorem --- and again in M11 with the
constraint cone of a linear program, where the result is LP duality. \emph{If those three
proofs do not look like one proof when M11 is reached, M1 was not finished.}
\item The \emph{polar cone} and \emph{Moreau's decomposition} of S5 are the normal cone M4 uses
to write optimality conditions, and the conic special case of the proximal decomposition proved
in M3.
\item The \emph{simplex projection} built in S3 is used verbatim by M8, M9 and M15; it is the
operation every projected policy-gradient step performs.
\item The \emph{closedness lemma} of S6 is why the polyhedral cone of M11 is closed, hence why
linear programming has no duality gap and needs no constraint qualification.
\end{itemize}

\section{Learning objectives}

On completing M1 the student should be able to do the following, without notes:

\begin{enumerate}
\item State and prove Carath\'eodory's theorem and its conic version, and say where finite
dimension is used.
\item Distinguish the interior from the relative interior, exhibit a convex set with empty
interior and nonempty relative interior, and say why every constraint qualification in the
plan is stated with $\ri$.
\item State and prove the projection theorem: existence, uniqueness, the variational
characterization, nonexpansiveness and firm nonexpansiveness.
\item Derive the projection onto a halfspace, an affine subspace, a ball, a box and the
probability simplex, and explain why clipping and renormalizing is \emph{not} the projection
onto the simplex.
\item State the three separation theorems --- proper, strong, supporting --- with their three
different hypotheses, and prove strong separation from the projection theorem in two lines,
exhibiting the separating vector explicitly.
\item Give two disjoint closed convex sets admitting no strong separation, and name the
hypothesis that fails.
\item Define the polar cone, state and prove the bipolar theorem, and prove Moreau's
decomposition, identifying the orthogonal decomposition of a subspace as its special case.
\item Prove that a finitely generated cone is closed, and exhibit a convex cone that is not.
\item State and prove Farkas' lemma in both its equality and inequality forms, marking the
single line where closedness is used, and derive Gordan's theorem from it.
\item Say, in one sentence for each of convexity, closedness and finite generation, which
theorem it enables and what breaks without it.
\end{enumerate}

\section{How to use this document}

\mpart{This document is self-contained.} Every definition it uses is stated in it, every
theorem it uses is proved in it, and no proof is left as ``see the reference''. There are
exactly two exceptions, both marked in place: Proposition \ref{prop:ricl}(2) and Theorem
\ref{thm:proper}, which are stated without proof because they are needed only for orientation
and their proofs are long and unrewarding at this stage. The books listed in
Table \ref{tab:sources} are for a \emph{second view} after a session, not preparation before
it. You should be able to work through M1 with this document, paper, and a Python interpreter.

\mpart{The order within a session is deliberate.} Following the expository method of Zhao's
\emph{Mathematical Foundations of Reinforcement Learning}, each session begins with the
question it answers and an informal picture of the answer, states the result next, and only
then proves it. Read in that order. A proof read before the statement it proves has been
understood as a claim about a picture is a proof that will not be remembered.

Each of the seven sections below is one working session, and each has these parts, always in
this order:

\begin{description}[leftmargin=1.9cm,style=nextline,font=\normalfont\itshape]
\item[The question] one paragraph. What is not yet known, stated so that the session's theorem
is recognizably its answer.
\item[The picture] the answer in words and geometry, before any formalism. If you cannot draw
this, do not proceed to the statements.
\item[Statements and proofs] definitions, theorems and complete proofs. These are to be read
once, then reproduced on paper with the document closed.
\item[Worked example] one instance carried through with explicit numbers, so that the theorem
becomes something you have computed rather than something you have read.
\item[Where this is used later] the forward links, named by module.
\item[Proof of concept] one implementation, sized to the session, with numbered validation
checks. Every check is an inequality against a stated tolerance or an exact identity. A check
reported without a number has not been performed.
\item[Exercises] written work, tagged where a later module reuses it.
\item[Session gate] the condition for starting the next session.
\end{description}

\mpart{The gate rule.} Do not start session $k+1$ until the gate of session $k$ has been
passed. The gates are cumulative by design: S7's capstone runs every earlier session's code
again, and if an earlier check was skipped it will fail there instead.

\mpart{Programming rules, once, for all sessions.} Python with NumPy and SciPy. Every object in
this module is one you construct, so every target is known in closed form or by an independent
solve; nothing is estimated from data. Fix and report one seed for the whole module. Geometric
identities are \emph{exact} and are checked at $10^{-12}$; comparisons against a numerical
solver are checked at that solver's stated tolerance and no better --- a check claiming
$10^{-12}$ agreement with an interior point solver is a check that was not performed. Where a
tolerance below is $10^{-9}$ rather than $10^{-12}$, the looser number is deliberate and marks
a quantity produced by an iterative solve. Keep one file per session, \texttt{s1.py} to
\texttt{s7.py}, plus a shared \texttt{m1lib.py} that S7 imports.

\mpart{One design decision, made now.} Represent every convex set by a small object exposing
three methods: \texttt{contains(x)}, \texttt{sample()} and \texttt{project(x)}. Represent a
cone additionally by generators or by an inequality description, with a converter between them.
Every validation check in S3 through S7 is then a loop over the same list of sets, and S4, S5
and S6 each cost a few lines instead of a rewrite. \emph{This is the single decision that
determines how long the module takes.}

\clearpage

\section{Session map}

\begin{table}[!ht]
\centering
\small
\caption{The seven sessions of M1, with the hour split and the deliverable of each. St, Wr and
Cd are hours of studying this text, writing out proofs and exercises by hand, and code.}
\label{tab:map}
\begin{tabular}{@{}l L{4.6cm} C{0.7cm}C{0.7cm}C{0.7cm}C{0.7cm} L{4.9cm}@{}}
\toprule
\textbf{S} & \textbf{Title} & \textbf{St} & \textbf{Wr} & \textbf{Cd} & \textbf{Tot}
& \textbf{Proof of concept} \\
\midrule
S1 & Convex sets, hulls and Carath\'eodory              & 0.3 & 0.3 & 0.2 & 0.8
   & The Carath\'eodory reduction \\
S2 & Relative interior, closure and closedness          & 0.3 & 0.3 & 0.2 & 0.8
   & The closedness failure, numerically \\
S3 & Projection onto a closed convex set                & 0.4 & 0.3 & 0.5 & 1.2
   & The projection library \\
S4 & Separation                                         & 0.3 & 0.5 & 0.2 & 1.0
   & \texttt{separate(C,x)} from \texttt{project} \\
S5 & Cones, polars and Moreau's decomposition           & 0.3 & 0.3 & 0.2 & 0.8
   & Polars and the conic decomposition \\
S6 & Closed cones and Farkas' lemma                     & 0.5 & 0.5 & 0.4 & 1.4
   & The Farkas certificate \\
S7 & Consolidation, capstone and checkpoint             & 0.2 & 0.3 & 0.5 & 1.0
   & The suite, plus the non-closed cone \\
\midrule
   & \textbf{Total} & \textbf{2.3} & \textbf{2.5} & \textbf{2.2} & \textbf{7.0} & \\
\bottomrule
\end{tabular}
\end{table}

\begin{table}[!ht]
\centering
\small
\caption{Sources for a second view. \emph{None of these is required.} Every proof this module
uses is written out below; the column says what each book adds if you want another treatment
after finishing a session.}
\label{tab:sources}
\begin{tabular}{@{}L{5.2cm} L{1.9cm} L{7.8cm}@{}}
\toprule
\textbf{Source} & \textbf{After} & \textbf{What it adds} \\
\midrule
Hiriart-Urruty and Lemar\'echal, \emph{Fundamentals of Convex Analysis}, part A
  & S1--S6
  & The same material with more examples and the infinite-dimensional caveats; A.3 and A.4 are
    the closest thing in print to S3 and S4 \\
Boyd and Vandenberghe, ch.~2, \S 8.1
  & S1, S3--S5
  & A large catalogue of convex sets and of closed-form projections; the dual-cone conventions
    used in the engineering literature \\
Crouzeix and Oca\~na, \emph{An\'alisis convexo}
  & S1--S6
  & The MMC007 syllabus text, in Spanish; useful for fixing vocabulary in the language of the
    lectures \\
Rockafellar, \emph{Convex Analysis}
  & S2, S4--S6
  & The reference statements later modules cite: \S 6, \S 11, \S 14, \S 17, \S 19, \S 22.
    Also the proofs omitted here (\S 6 and Thm.~11.3) \\
Bertsekas, \emph{Nonlinear Programming}, App.~B
  & S2, S6
  & Recession cones in far more detail than S2 needs, and the closedness criteria for images
    and sums \\
Schrijver, \S 7.3
  & S6
  & Farkas' variants collected, with the polyhedral rather than the separation proof --- worth
    seeing once, as a contrast \\
Bertsimas and Tsitsiklis, \S 4.6
  & S6
  & Farkas read as an infeasibility certificate; this is the form M11 and M12 use \\
Beck, \emph{First-Order Methods}, ch.~6
  & S3
  & A larger projection table, and the projection as a special proximal operator, which is M3 \\
Farkas (1902); Gordan (1873); Moreau (1962)
  & S5, S6
  & The originals of Theorems \ref{thm:farkas}, \ref{thm:gordan} and \ref{thm:moreau}; Moreau
    is two pages \\
\bottomrule
\end{tabular}
\end{table}

\clearpage

\section{Setup and notation}

The ambient space is $\R^n$ with the Euclidean inner product $\ip{\cdot}{\cdot}$ and the norm
$\norm{x}=\ip{x}{x}^{1/2}$. Everything here is finite-dimensional; the two places where that is
essential rather than convenient are marked when they occur. Convex \emph{functions} belong to
M2 and are not used, with one exception flagged in S3.

\begin{definition}[Convex set, hulls, cone]
\label{def:convex}
$C\subseteq\R^n$ is \emph{convex} if $(1-\lambda)x+\lambda y\in C$ whenever $x,y\in C$ and
$\lambda\in[0,1]$. A \emph{convex combination} of $x_1,\dots,x_k$ is $\sum\lambda_ix_i$ with
$\lambda_i\ge0$, $\sum\lambda_i=1$; an \emph{affine} combination drops $\lambda_i\ge0$; a
\emph{conic} combination drops $\sum\lambda_i=1$. $\conv S$, $\aff S$ and $\cone S$ denote the
intersections of all convex, affine and convex-conic sets containing $S$. A set $K$ is a
\emph{cone} if $\lambda x\in K$ for all $x\in K,\lambda\ge0$; note $0\in K$ for every cone here.
\end{definition}

\begin{definition}[Relative interior]
\label{def:ri}
$\ri C$ is the interior of $C$ relative to $\aff C$: the set of $x\in C$ for which there is
$\varepsilon>0$ with $B(x,\varepsilon)\cap\aff C\subseteq C$. The relative boundary is
$\cl C\setminus\ri C$.
\end{definition}

\begin{definition}[Projection and distance]
\label{def:proj}
For $C$ nonempty closed convex, $\proj_Cx=\argmin_{c\in C}\norm{x-c}$ and
$\dist_C(x)=\norm{x-\proj_Cx}$. Theorem \ref{thm:proj} shows both are well defined.
\end{definition}

\begin{definition}[Polar and dual cone]
\label{def:polar}
For $S\subseteq\R^n$, the \emph{polar} is $S^\circ=\{y:\ip{y}{s}\le0\ \forall s\in S\}$ and the
\emph{dual} is $S^*=-S^\circ=\{y:\ip{y}{s}\ge0\ \forall s\in S\}$. \emph{This document uses the
polar throughout.}
\end{definition}

Two abbreviations used without comment: $\Simplex_n=\{x\in\R^n:x\ge0,\ \sum_ix_i=1\}$ is the
probability simplex, and for $A\in\R^{m\times n}$, $\cone(A)=\{Ax:x\ge0\}$ is the cone generated
by the columns of $A$.

Four warnings, all of which cause avoidable confusion later:

\begin{itemize}
\item \textbf{``Interior'' is the wrong word.} The feasible set $\{x:Ax=b,\ x\ge0\}$ of a linear
program in standard form has empty interior in $\R^n$ whenever $A$ has a nonzero row, and so
does $\Simplex_n$. Every theorem in this plan that appears to require an interior point ---
Slater's condition in M4, the strictly feasible start of M13 --- is stated with the relative
interior of Definition \ref{def:ri}. S2 exists to make that distinction reflexive.

\item \textbf{Closedness is a hypothesis, not a technicality.} Strong separation is false
without it; the sum of two closed convex sets need not be closed; and Farkas' lemma is an
\emph{alternative} --- exactly one of two systems holds --- rather than a mere implication only
because $\cone(A)$ is closed, which is Lemma \ref{lem:closedcone} and is the one genuinely
nontrivial lemma of the module. The capstone in S7 exhibits what happens when it fails, and the
answer is a nonzero duality gap.

\item \textbf{$\proj_C$ is not linear.} It is nonexpansive, which is weaker, and it is linear
exactly when $C$ is an affine subspace --- the case MMC201's module M2 studies under the name
``least squares''. Expecting $\proj_C$ to commute with a gradient step is the standard first
error in reading the projected gradient method of M8, and S3 check 5 exists to make the
non-linearity a number you have seen.

\item \textbf{Three separations, three hypotheses.} \emph{Proper} separation, \emph{strong}
separation (with a positive gap), and a \emph{supporting} hyperplane (through a boundary point)
are three different theorems. Which one is available decides whether the duality gap of M4 is
zero. Do not use the word ``separating'' without saying which.
\end{itemize}

\clearpage

\section{Session S1. Convex sets, hulls and Carath\'eodory}

\mpart{The question.} The convex hull $\conv S$ is defined as an intersection over all convex
sets containing $S$ --- an uncountable family --- or, equivalently, as the set of \emph{all}
convex combinations of \emph{finitely many} points of $S$, with no bound on how many. Neither
description can be checked by a computer. If somebody claims $x\in\conv S$, what is the smallest
amount of evidence that settles it?

\mpart{The picture.} Take a cloud of fifty points in the plane and a point $x$ inside their
hull. You never need fifty of them: find a triangle with vertices in the cloud that contains
$x$, and three points suffice. The claim of this session is that ``three'' is ``$n+1$'' in
$\R^n$, and that this is a dimension count rather than an accident. Here is the mechanism. If
$x$ is written using more than $n+1$ points, those points are affinely dependent. An affine
dependence is a direction in weight-space along which the weights can be moved \emph{without
moving $x$}. Slide along it until some weight first hits zero. One point has been eliminated
and $x$ has not moved. Repeat.

That last paragraph is the entire proof. What follows writes it down.

\subsection*{Statements and proofs}

\begin{lemma}[The hull is the set of finite convex combinations]
\label{lem:hull}
$\conv S$ equals the set $T$ of all convex combinations of finitely many points of $S$.
\end{lemma}

\begin{proof}
$T$ is convex: a convex combination of two convex combinations is a convex combination. $T$
contains $S$ (take $k=1$). Hence $\conv S\subseteq T$.

For $T\subseteq\conv S$, let $C$ be any convex set with $S\subseteq C$; we show $T\subseteq C$
by induction on the number $k$ of atoms. For $k=1$ this is $S\subseteq C$. Suppose it holds for
$k-1$ and let $x=\sum_{i=1}^k\lambda_ix_i$ with $x_i\in S$, $\lambda_i\ge0$, $\sum\lambda_i=1$.
If $\lambda_k=1$ then $x=x_k\in C$. Otherwise $\lambda_k<1$, and
\[
x=(1-\lambda_k)\,y+\lambda_kx_k,\qquad
y=\sum_{i=1}^{k-1}\frac{\lambda_i}{1-\lambda_k}\,x_i .
\]
The coefficients of $y$ are nonnegative and sum to $1$, so $y\in C$ by the inductive hypothesis,
and then $x\in C$ by convexity. Since $C$ was arbitrary, $T\subseteq\conv S$.
\end{proof}

\begin{theorem}[Carath\'eodory]
\label{thm:cara}
Let $S\subseteq\R^n$ and $x\in\conv S$. Then $x$ is a convex combination of at most $n+1$
points of $S$.
\end{theorem}

\begin{proof}
By Lemma \ref{lem:hull}, write $x=\sum_{i=1}^k\lambda_ix_i$ with $x_i\in S$, all
$\lambda_i>0$ (discard zero weights), $\sum\lambda_i=1$, and $k$ \emph{minimal} among all such
representations. Suppose, for contradiction, that $k\ge n+2$.

The $k-1\ge n+1$ vectors $x_2-x_1,\dots,x_k-x_1$ lie in $\R^n$ and so are linearly dependent:
there are $\mu_2,\dots,\mu_k$, not all zero, with $\sum_{i=2}^k\mu_i(x_i-x_1)=0$. Put
$\mu_1=-\sum_{i=2}^k\mu_i$. Then
\[
\sum_{i=1}^k\mu_ix_i=0,\qquad \sum_{i=1}^k\mu_i=0,\qquad \mu\ne0 .
\]
Because the $\mu_i$ sum to zero and are not all zero, at least one is strictly positive. Let
\[
t^\star=\min\Big\{\frac{\lambda_i}{\mu_i}\ :\ \mu_i>0\Big\},
\]
attained at some index $j$, and set $\lambda_i'=\lambda_i-t^\star\mu_i$. Then:
\begin{itemize}
\item $\sum_i\lambda_i'=\sum_i\lambda_i-t^\star\sum_i\mu_i=1$;
\item $\sum_i\lambda_i'x_i=x-t^\star\cdot0=x$;
\item $\lambda_i'\ge0$: if $\mu_i>0$ then $\lambda_i'=\mu_i(\lambda_i/\mu_i-t^\star)\ge0$ by the
definition of $t^\star$; if $\mu_i\le0$ then $\lambda_i'\ge\lambda_i>0$ since $t^\star\ge0$;
\item $\lambda_j'=0$ exactly.
\end{itemize}
So $x$ is a convex combination of at most $k-1$ points of $S$, contradicting the minimality of
$k$. Hence $k\le n+1$.
\end{proof}

\begin{remark}
Finite dimension is used once, in ``$k-1\ge n+1$ vectors in $\R^n$ are linearly dependent''.
In an infinite-dimensional space the theorem is false as stated; the correct substitute is a
statement about closures, and it is one of the reasons infinite-dimensional convex analysis is
a different subject.
\end{remark}

\begin{corollary}[Conic Carath\'eodory]
\label{cor:conecara}
If $x\in\cone S$ with $S\subseteq\R^n$, then $x=\sum_{i=1}^k\mu_is_i$ with $\mu_i>0$, $s_i\in S$
and $s_1,\dots,s_k$ \emph{linearly independent}; in particular $k\le n$.
\end{corollary}

\begin{proof}
Take such a representation with $k$ minimal and all $\mu_i>0$. If $s_1,\dots,s_k$ were linearly
dependent there would be $\nu\ne0$ with $\sum_i\nu_is_i=0$; replacing $\nu$ by $-\nu$ if
necessary, some $\nu_i>0$. With $t^\star=\min\{\mu_i/\nu_i:\nu_i>0\}$ attained at $j$, the
coefficients $\mu_i-t^\star\nu_i$ are nonnegative, represent $x$, and vanish at $j$,
contradicting minimality. So the $s_i$ are independent and $k\le n$.
\end{proof}

\begin{proposition}
\label{prop:convcompact}
If $K\subseteq\R^n$ is compact then $\conv K$ is compact.
\end{proposition}

\begin{proof}
Let $\Phi:\Simplex_{n+1}\times K^{n+1}\to\R^n$, $\Phi(\lambda,x_1,\dots,x_{n+1})
=\sum_{i=1}^{n+1}\lambda_ix_i$. It is continuous and its domain is compact, so its image is
compact. The image is contained in $\conv K$; and it contains $\conv K$ by Theorem
\ref{thm:cara}, since a combination of $k\le n+1$ points becomes one of exactly $n+1$ points by
repeating a point with weight $0$. Hence $\conv K$ is that image.
\end{proof}

Compactness cannot be weakened to closedness, and the following example is worth keeping: S2
explains what has to be added, S4 uses a relative of it, and S7 uses it again.

\begin{example}[The convex hull of a closed set need not be closed]
\label{ex:hullnotclosed}
Let $S=\{(u,v):u>0,\ v\ge1/u\}\cup\{(0,0)\}\subseteq\R^2$. The first piece is closed (a
sequence in it with $u_k\to0^+$ has $v_k\to\infty$, so has no limit in $\R^2$), the second is a
point, and a union of two closed sets is closed; so $S$ is closed. But $\conv S$ is not.
Indeed, for $\varepsilon\in(0,1]$ take $u=\varepsilon^{-1/2}$ and the combination
\[
(1-\sqrt\varepsilon)\,(0,0)+\sqrt\varepsilon\,\big(\varepsilon^{-1/2},\ \varepsilon^{1/2}\big)
=(1,\varepsilon)\in\conv S ,
\]
so $(1,\varepsilon)\in\conv S$ for every $\varepsilon>0$. Yet $(1,0)\notin\conv S$: every point
of $S$ other than the origin has $v>0$, so any convex combination with $v=0$ must place all its
weight on the origin and is therefore $(0,0)$. Thus $(1,\varepsilon)\to(1,0)\notin\conv S$.
\end{example}

\subsection*{Worked example}

In $\R^2$ take $S=\{(1,0),(0,1),(-1,0),(0,-1)\}$ and $x=(0,0)=\tfrac14\sum_{i}s_i$. This uses
$k=4>n+1=3$ atoms, so Theorem \ref{thm:cara} says it can be reduced. Run the proof.

The four points are affinely dependent: $(1,0)-(0,1)+(-1,0)-(0,-1)=0$ with coefficients
$\mu=(1,-1,1,-1)$ summing to $0$. The positive entries are $\mu_1=\mu_3=1$, and
\[
t^\star=\min\{\lambda_1/\mu_1,\ \lambda_3/\mu_3\}=\min\{\tfrac14,\tfrac14\}=\tfrac14 .
\]
Then $\lambda'=\lambda-t^\star\mu=(0,\tfrac12,0,\tfrac12)$, giving
$x=\tfrac12(0,1)+\tfrac12(0,-1)$: two atoms, and indeed $2\le3$. Note that \emph{two} weights
vanished at once here, because the minimum in $t^\star$ was attained twice; the proof only
promises one.

\subsection*{Where this is used later}

Corollary \ref{cor:conecara} is the engine of Lemma \ref{lem:closedcone} in S6, which is the
hypothesis of Farkas' lemma, which is the theorem the rest of the plan quotes most. The
module's logical chain runs through its shortest session. Theorem \ref{thm:cara} also bounds
the support of a Frank--Wolfe iterate in M8 and explains the vertex solutions of M11.

\mpart{Reinforcement learning.} A stochastic policy at one state is a point of
$\Simplex_{|\mathcal A|}=\conv\{\delta_a\}$, so Theorem \ref{thm:cara} there says only that a
distribution is a mixture of point masses. The content appears one level up. The set of
discounted occupancy measures of a finite Markov decision process is a polytope in
$\R^{|\mathcal S||\mathcal A|}$ whose vertices are the deterministic policies, so \emph{every
achievable occupancy measure --- hence every achievable behaviour --- is a mixture of at most
$|\mathcal S||\mathcal A|+1$ deterministic policies.} That finiteness is why M11's linear
program has a vertex solution and why one may search over deterministic policies without loss.

\subsection*{Proof of concept: the Carath\'eodory reduction}

Implement \texttt{caratheodory(x, S)} by iterating the elimination step of the proof: solve for
an affine dependence (a null vector of the matrix with rows $x_i^\top$ and a row of ones),
compute $t^\star$, update, drop the vanished atom. Implement the set objects for the box, the
ball, the simplex and a random polytope, each with \texttt{contains} and \texttt{sample};
\texttt{project} arrives in S3.

\emph{Validation.}
\begin{enumerate}
\item On $10^3$ random instances with $n\in\{2,5,10\}$ and $|S|=50$: the output has at most
$n+1$ atoms, weights are nonnegative and sum to $1$ within $10^{-12}$, and the combination
reconstructs $x$ within $10^{-12}$.
\item The reduction is exact, not approximate: assert that at each step the eliminated weight
is \emph{exactly} $0$, by construction rather than by tolerance. If this fails, $t^\star$ was
taken as a minimum over the wrong index set --- the minimum runs over $\mu_i>0$ only.
\item Reproduce the worked example above and confirm the routine returns two atoms with weights
$\tfrac12,\tfrac12$, not three.
\item Feed a point already supported on $k\le n+1$ atoms and confirm the routine returns it
unchanged. This catches the common bug of eliminating unconditionally.
\end{enumerate}

\emph{Deliverable.} \texttt{s1.py}, and \texttt{caratheodory} plus the set objects in
\texttt{m1lib.py}.

\subsection*{Exercises}
\begin{enumerate}[label=E1.\arabic*]
\item Reproduce the proof of Theorem \ref{thm:cara} with the document closed. Mark the one
sentence that uses finite dimension.
\item Show that an arbitrary intersection of convex sets is convex, and give two convex sets
whose union is not. Deduce that the feasible set of an optimization problem is convex as soon
as each constraint set is --- and note that this closure property is the entire practical
reason convex programming is a usable class of problems.
\item Verify every claim in Example \ref{ex:hullnotclosed}, including the arithmetic of the
displayed combination.
\item Prove Corollary \ref{cor:conecara} without looking, and show
$\cone\{a_1,\dots,a_m\}=\{Ax:x\ge0\}$ where $A$ has the $a_i$ as columns. \textbf{[Reused in
S6: this corollary is the proof of Lemma \ref{lem:closedcone}.]}
\end{enumerate}

\mpart{Session gate.} Checks 1--4 pass and E1.4 is written out.

\clearpage

\section{Session S2. Relative interior, closure and the closedness hypothesis}

\mpart{The question.} Every theorem in the rest of this module says ``closed convex''. It is
tempting to read ``closed'' as a technical courtesy. Is it? And separately: optimization
constantly produces sets with \emph{empty interior} --- anything cut out by an equality
constraint --- so what replaces ``interior point'' in the hypotheses of the theorems that seem
to need one?

\mpart{The picture.} Two pictures, one for each half of the question.

First: the probability simplex in $\R^3$ is a solid triangle sitting inside a plane. It has no
interior point in $\R^3$ --- every ball around any of its points sticks out of the plane. But
inside the plane it plainly has an inside. The fix is to measure interiority relative to the
affine hull, and it is not a cosmetic fix: the entire hypothesis of the strong duality theorem
in M4 is of this form.

Second: take the region above a hyperbola in the right half-plane, and the vertical axis. Both
are closed and convex. Add them. Going far to the right along the hyperbola, the region's points
come arbitrarily close to the horizontal axis; sliding those points vertically by any amount ---
which is what adding the axis does --- sweeps out the entire \emph{open} right half-plane. The
sum has lost its boundary, because the two sets escape to infinity in complementary directions
and the sum accumulates on points that neither sum ever reaches. This is not pathology. It is
the mechanism of every duality gap in the plan.

\subsection*{Statements and proofs}

\begin{proposition}
\label{prop:rinonempty}
If $C\subseteq\R^n$ is convex and nonempty then $\ri C\ne\emptyset$.
\end{proposition}

\begin{proof}
Let $d=\dim\aff C$ and choose $x_0,\dots,x_d\in C$ affinely independent (possible by the
definition of $d$). The map $\psi(\lambda_1,\dots,\lambda_d)=x_0+\sum_{i=1}^d\lambda_i(x_i-x_0)$
is an affine bijection from $\R^d$ onto $\aff C$, and it is a homeomorphism because the
$x_i-x_0$ are linearly independent. The point
$b=\psi(\tfrac1{d+1},\dots,\tfrac1{d+1})=\tfrac1{d+1}\sum_{i=0}^dx_i$ is the barycentre of the
simplex $\conv\{x_0,\dots,x_d\}\subseteq C$. Its preimage lies in the open set
$\{\lambda:\lambda_i>0,\ \sum\lambda_i<1\}$, whose image under $\psi$ consists of convex
combinations of the $x_i$ with all weights positive, hence lies in $C$. As $\psi$ is a
homeomorphism onto $\aff C$, that image is relatively open, so $b\in\ri C$.
\end{proof}

\begin{proposition}[Line segment principle]
\label{prop:segment}
If $x\in\ri C$, $y\in\cl C$ and $\lambda\in[0,1)$, then $z=(1-\lambda)x+\lambda y\in\ri C$.
\end{proposition}

\begin{proof}
Work inside $\aff C$, so we may assume $\aff C=\R^n$ and $\ri=\mathrm{int}$. Let $\varepsilon>0$
with $B(x,\varepsilon)\subseteq C$, and set $r=(1-\lambda)\varepsilon/2>0$. Let $w$ satisfy
$\norm{w-z}<r$. Since $y\in\cl C$ there is $y'\in C$ with $\norm{y-y'}<r$. Define
\[
x'=\frac{w-\lambda y'}{1-\lambda},\qquad\text{so that}\qquad w=(1-\lambda)x'+\lambda y' .
\]
Then, using $z=(1-\lambda)x+\lambda y$,
\[
\norm{x'-x}=\frac{\norm{w-\lambda y'-(1-\lambda)x}}{1-\lambda}
=\frac{\norm{(w-z)+\lambda(y-y')}}{1-\lambda}
\le\frac{r+\lambda r}{1-\lambda}\le\frac{2r}{1-\lambda}=\varepsilon .
\]
Hence $x'\in B(x,\varepsilon)\subseteq C$, and $w=(1-\lambda)x'+\lambda y'\in C$ by convexity.
So $B(z,r)\subseteq C$ and $z\in\ri C$.
\end{proof}

\begin{proposition}
\label{prop:ricl}
For $C$ convex nonempty: (1) $\cl(\ri C)=\cl C$; (2) $\ri(\cl C)=\ri C$.
\end{proposition}

\begin{proof}
(1) $\subseteq$ is immediate. For $\supseteq$, take $y\in\cl C$ and, by Proposition
\ref{prop:rinonempty}, some $x\in\ri C$. By Proposition \ref{prop:segment},
$(1-\lambda)x+\lambda y\in\ri C$ for every $\lambda<1$; letting $\lambda\to1^-$ exhibits $y$ as
a limit of points of $\ri C$. Statement (2) is proved similarly and is stated here without proof
--- it is used only for orientation, and its proof adds nothing this module needs.
\end{proof}

\begin{example}[The simplex]
\label{ex:simplexri}
$\Simplex_n\subseteq H=\{x:\sum_ix_i=1\}$, a hyperplane, which has empty interior in $\R^n$;
hence $\mathrm{int}\,\Simplex_n=\emptyset$. On the other hand $\aff\Simplex_n=H$, and
$\ri\Simplex_n=\{x\in H:x_i>0\ \forall i\}$: if all $x_i>0$ then any $y\in H$ with
$\norm{y-x}<\min_ix_i$ has $y_i>0$, so $y\in\Simplex_n$; and if some $x_j=0$, then
$x+\delta(e_j-e_k)\in H$ for $k\ne j$ has a negative $j$-th coordinate for every $\delta<0$
arbitrarily small in norm, so $x\notin\ri\Simplex_n$.
\end{example}

\begin{example}[A sum of closed convex sets that is not closed]
\label{ex:notclosed}
Let $C=\{(u,v):u>0,\ v\ge1/u\}$ and $D=\{0\}\times\R$. Both are convex, and both are closed
($C$ by the argument of Example \ref{ex:hullnotclosed}). Their sum is
\[
C+D=\{(u,v+t):u>0,\ v\ge1/u,\ t\in\R\}=\{(u,s):u>0,\ s\in\R\},
\]
the \emph{open} right half-plane, which is not closed: $(1/k,0)\in C+D$ for every $k$, and the
limit $(0,0)$ is not in $C+D$. The same set shows that the affine image of a closed convex set
need not be closed, since the projection of $C$ onto its first coordinate is $(0,\infty)$.
\end{example}

Why does this happen, and when does it not? The answer is about directions of escape.

\begin{definition}[Recession cone]
\label{def:recession}
For $C$ closed convex nonempty, $0^+C=\{d:x+td\in C\ \forall x\in C,\ t\ge0\}$.
\end{definition}

\begin{lemma}
\label{lem:reclimit}
Let $C$ be closed convex, $c_k\in C$ with $\norm{c_k}\to\infty$ and $c_k/\norm{c_k}\to u$. Then
$u\in0^+C$.
\end{lemma}

\begin{proof}
Fix $x\in C$ and $t\ge0$. For $k$ large, $t/\norm{c_k}\le1$, so by convexity
$x+\tfrac{t}{\norm{c_k}}(c_k-x)\in C$. As $k\to\infty$ this converges to $x+tu$, and $C$ is
closed, so $x+tu\in C$.
\end{proof}

\begin{proposition}[Recession criterion]
\label{prop:recession}
Let $C,D$ be closed convex nonempty with $0^+C\cap(-0^+D)=\{0\}$ --- in particular, with either
set bounded. Then $C+D$ is closed.
\end{proposition}

\begin{proof}
Let $z_k=c_k+d_k\to z$ with $c_k\in C$, $d_k\in D$. If $(c_k)$ has a bounded subsequence, pass
to it and to a convergent sub-subsequence $c_k\to c\in C$ (closed); then $d_k=z_k-c_k\to z-c$,
which lies in $D$ (closed), so $z\in C+D$. Otherwise $\norm{c_k}\to\infty$, and since $(z_k)$ is
bounded, $\norm{d_k}\to\infty$ with $\norm{d_k}/\norm{c_k}\to1$. Pass to a subsequence with
$c_k/\norm{c_k}\to u$, $\norm u=1$; then
$d_k/\norm{d_k}=\tfrac{\norm{c_k}}{\norm{d_k}}\cdot\tfrac{z_k-c_k}{\norm{c_k}}\to-u$. By Lemma
\ref{lem:reclimit}, $u\in0^+C$ and $-u\in0^+D$, so $u\in0^+C\cap(-0^+D)=\{0\}$, contradicting
$\norm u=1$. (If either set is bounded its recession cone is $\{0\}$ and the hypothesis holds
automatically.)
\end{proof}

In Example \ref{ex:notclosed}, $0^+C=\R_+^2$ and $0^+D=-0^+D=\{0\}\times\R$, so
$0^+C\cap(-0^+D)=\{0\}\times\R_+\ne\{0\}$: the hypothesis fails exactly where the closedness
does, and the offending direction $(0,1)$ is visibly the direction in which the two sets escape.

\subsection*{Worked example}

Take $C$ and $D$ of Example \ref{ex:notclosed} and ask for the pair $(c,d)$ minimizing
$\norm{c+d}$. Parameterize $c=(u,1/u)$ with $u>0$ and $d=(0,t)$; then
$\norm{c+d}^2=u^2+(1/u+t)^2$, and for each $u$ the best $t$ is $t=-1/u$, giving $\norm{c+d}=u$.
So the infimum over the parameterization is $\inf_{u>0}u=0$, \emph{and it is attained at no
$u>0$}: the minimizing sequence is $u_k\to0^+$, at which $\norm{d_k}=1/u_k\to\infty$. The
infimum $0$ is the distance from the origin to $C+D$, the origin is in $\cl(C+D)$, and the
minimizer has escaped to infinity in the direction $(0,-1)\in-0^+D$ that Proposition
\ref{prop:recession} forbids. \emph{An infimum approached only at infinity is what ``not
closed'' means operationally}, and this is the phrase to remember.

Now intersect $C$ with the ball $\bar B(0,R)$, which makes it compact. The constraint
$u^2+1/u^2\le R^2$ forces $u\ge u_{\min}(R)>0$, so the minimum becomes $u_{\min}(R)>0$ and is
attained on the sphere. Compactness restores attainment and the optimal value is strictly
positive for every finite $R$, which is Proposition \ref{prop:recession} seen from the other
side.

\subsection*{Where this is used later}

Every ``closed convex'' in S3--S6 is this session. Concretely: Theorem \ref{thm:proj} needs
closedness for existence; Theorem \ref{thm:septwo} needs Proposition \ref{prop:recession} to
know $C-D$ is closed; Lemma \ref{lem:closedcone} is a closedness statement and is the hypothesis
of Farkas. In M4, the duality gap is nonzero exactly when a certain perturbation set fails to be
closed, and the mechanism is the one in Example \ref{ex:notclosed}.

\mpart{Reinforcement learning.} The set of discounted occupancy measures of a finite Markov
decision process is a polytope, hence compact, so by Proposition \ref{prop:convcompact} and
Weierstrass the supremum of expected return over policies is attained. M11 may therefore write
$\max$ rather than $\sup$, and an optimal policy exists --- a geometric fact, not a
probabilistic one.

\subsection*{Proof of concept: the closedness failure, numerically}

Implement the parameterized minimization of the worked example on a grid over $u$, refined five
times, with the optimal $t=-1/u$ substituted in closed form.

\emph{Validation.}
\begin{enumerate}
\item Tabulate, at five successive refinements, the best value $\norm{c+d}$ and the minimizing
$u$ and $\norm d=1/u$. Confirm the value decreases toward $0$ while $\norm{d}$ roughly doubles
at each halving of the grid tolerance. \emph{Print the table.}
\item Repeat with $C\cap\bar B(0,R)$ for $R\in\{2,10,100\}$: confirm the minimum is now attained
at an interior grid point, that the minimizer satisfies $\norm c=R$ to the grid resolution, and
that the optimal value is strictly positive and decreasing in $R$. \emph{This tests Proposition
\ref{prop:recession} as an if-and-only-if rather than illustrating it.}
\item Example \ref{ex:simplexri} numerically for $n=3$: confirm by sampling that no point of
$\Simplex_3$ has a full-dimensional ball inside it --- every random perturbation in $\R^3$
leaves the affine hull --- while for each sampled $x$ with all $x_i>0$ you can exhibit an
explicit $\varepsilon$ such that every $y$ in $\aff\Simplex_3$ with $\norm{y-x}<\varepsilon$
lies in $\Simplex_3$. Report the empirical relation between $\varepsilon$ and $\min_ix_i$;
the proof of Example \ref{ex:simplexri} predicts $\varepsilon=\min_ix_i$.
\end{enumerate}

\emph{Deliverable.} \texttt{s2.py} and the divergence table of check 1.

\subsection*{Exercises}
\begin{enumerate}[label=E2.\arabic*]
\item Reproduce Proposition \ref{prop:segment} with the document closed, and deduce
Proposition \ref{prop:ricl}(1).
\item Verify Example \ref{ex:simplexri} in full, and state the consequence for how Slater's
condition must be written. \textbf{[Reused in M4.]}
\item Verify Example \ref{ex:notclosed}, compute both recession cones, and identify the
direction in which Proposition \ref{prop:recession} fails.
\item Prove that the image of a \emph{polyhedron} $\{x:Ax\le b\}$ under a linear map is a
polyhedron, hence closed, so that no counterexample of Example \ref{ex:notclosed}'s kind can be
polyhedral. The argument is Fourier--Motzkin elimination: eliminate one variable at a time by
pairing each inequality in which it has a positive coefficient with each in which it has a
negative one. \textbf{[Reused in M14: projecting out the complicating variables is exactly this
operation, and it is why Benders' master problem is a linear program.]}
\end{enumerate}

\mpart{Session gate.} Checks 1 and 2 pass; E2.2 and E2.4 are written out.

\clearpage

\section{Session S3. Projection onto a closed convex set}

\mpart{The question.} Given a point $x$ and a set $C$, is there a nearest point of $C$ to $x$?
Is it unique? And --- the question that actually matters --- is there a way to \emph{recognize}
the nearest point, some equation or inequality it satisfies and no other point of $C$ does?
Without such a characterization the nearest point is a definition and nothing follows from it.

\mpart{The picture.} For a subspace you know the answer: drop a perpendicular, and the
characterization is $x-p\perp C$. For a general convex set the perpendicular cannot be exact,
because $C$ has a boundary and $p$ may sit on it; the equality must weaken to an inequality.
Stand at $x$ and look at $p$, the nearest point. If some $y\in C$ were at an angle strictly less
than $90^\circ$ from the direction $x-p$, you could walk a little way from $p$ towards $y$ ---
staying in $C$, by convexity --- and get closer to $x$. So no $y\in C$ is at an angle below
$90^\circ$:
\[
\ip{x-p}{y-p}\le0\qquad\text{for every }y\in C .
\]
That inequality is the whole session. Everything else --- existence, uniqueness,
nonexpansiveness, every closed formula, the separating hyperplane of S4, Moreau's decomposition
in S5, the Farkas certificate of S6 --- is obtained by using it.

\subsection*{Statements and proofs}

\begin{theorem}[Projection theorem]
\label{thm:proj}
Let $C\subseteq\R^n$ be nonempty, closed and convex, and let $x,z\in\R^n$.
\begin{enumerate}
\item There is exactly one point $\proj_Cx\in C$ minimizing $\norm{x-c}$ over $c\in C$.
\item For $p\in C$: $p=\proj_Cx$ if and only if $\ip{x-p}{y-p}\le0$ for every $y\in C$.
\item $\norm{\proj_Cx-\proj_Cz}^2\le\ip{x-z}{\proj_Cx-\proj_Cz}$; consequently
$\norm{\proj_Cx-\proj_Cz}\le\norm{x-z}$.
\end{enumerate}
\end{theorem}

\begin{proof}
\emph{(1) Existence.} Fix any $c_0\in C$ and put $R=\norm{x-c_0}$. The set
$C\cap\bar B(x,R)$ is nonempty (it contains $c_0$), bounded, and closed --- \emph{this is the
only place closedness of $C$ is used} --- hence compact. The continuous function
$c\mapsto\norm{x-c}$ attains its minimum on it, at some $p$. Any $c\in C$ outside the ball has
$\norm{x-c}>R\ge\norm{x-p}$, so $p$ minimizes over all of $C$.

\emph{Uniqueness.} Let $p,q\in C$ both attain the minimum value $\delta$. Their midpoint
$m=(p+q)/2$ lies in $C$ by convexity --- \emph{this is where convexity is used} --- and the
parallelogram identity gives
\[
\norm{x-p}^2+\norm{x-q}^2=2\norm{x-m}^2+\tfrac12\norm{p-q}^2 ,
\]
that is, $2\delta^2=2\norm{x-m}^2+\tfrac12\norm{p-q}^2$. Since $\norm{x-m}\ge\delta$, we get
$\tfrac12\norm{p-q}^2\le0$, so $p=q$.

\emph{(2) ($\Leftarrow$).} Suppose $p\in C$ satisfies the inequality. For any $y\in C$,
\[
\norm{x-y}^2=\norm{(x-p)-(y-p)}^2=\norm{x-p}^2-2\ip{x-p}{y-p}+\norm{y-p}^2\ \ge\ \norm{x-p}^2 ,
\]
so $p$ is a minimizer, and by (1) it is \emph{the} minimizer.

\emph{($\Rightarrow$).} Let $p=\proj_Cx$ and $y\in C$. For $t\in(0,1]$ the point $p+t(y-p)$ lies
in $C$ by convexity, so
\[
\norm{x-p-t(y-p)}^2\ \ge\ \norm{x-p}^2
\quad\Longrightarrow\quad
-2t\ip{x-p}{y-p}+t^2\norm{y-p}^2\ \ge\ 0 .
\]
Divide by $t>0$ to get $-2\ip{x-p}{y-p}+t\norm{y-p}^2\ge0$, and let $t\downarrow0$:
$\ip{x-p}{y-p}\le0$.

\emph{(3).} Write $p=\proj_Cx$, $q=\proj_Cz$. Applying (2) at $x$ with $y=q$, and at $z$ with
$y=p$:
\[
\ip{x-p}{q-p}\le0,\qquad \ip{z-q}{p-q}\le0 .
\]
Adding them and writing $\ip{z-q}{p-q}=\ip{q-z}{q-p}$ gives $\ip{(x-z)-(p-q)}{q-p}\le0$, i.e.
$-\ip{x-z}{p-q}+\norm{p-q}^2\le0$, which is the claim. Cauchy--Schwarz then gives
$\norm{p-q}^2\le\norm{x-z}\norm{p-q}$, hence $\norm{p-q}\le\norm{x-z}$.
\end{proof}

\begin{remark}
Part (3) is called \emph{firm} nonexpansiveness and is strictly stronger than the
nonexpansiveness deduced from it. M9 needs the firm version: it is what makes proximal iterates
Fej\'er monotone with respect to the solution set, and plain nonexpansiveness is not enough for
that. Notice also that (3) was proved by using (2) twice and adding --- a manoeuvre that returns
in S5.
\end{remark}

\begin{proposition}
\label{prop:dist}
$\dist_C$ is convex and $1$-Lipschitz, and $\varphi=\tfrac12\dist_C^2$ is differentiable
everywhere with $\nabla\varphi(x)=x-\proj_Cx$.
\end{proposition}

\begin{proof}
\emph{Convexity.} For $x,y$ and $\lambda\in[0,1]$, the point
$\lambda\proj_Cx+(1-\lambda)\proj_Cy$ lies in $C$, so
\[
\dist_C(\lambda x+(1-\lambda)y)\le\norm{\lambda(x-\proj_Cx)+(1-\lambda)(y-\proj_Cy)}
\le\lambda\dist_C(x)+(1-\lambda)\dist_C(y).
\]
\emph{Lipschitz.} $\dist_C(x)\le\norm{x-\proj_Cz}\le\norm{x-z}+\dist_C(z)$, and symmetrically.

\emph{Differentiability.} Put $g=x-\proj_Cx$. Using $\proj_Cx\in C$ as a competitor for $y$,
\[
\varphi(y)\le\tfrac12\norm{y-\proj_Cx}^2=\varphi(x)+\ip{g}{y-x}+\tfrac12\norm{y-x}^2 ,
\]
which bounds the error above by $\tfrac12\norm{y-x}^2$. Using $\proj_Cy$ as a competitor at $x$,
\[
\varphi(x)\le\tfrac12\norm{x-\proj_Cy}^2=\varphi(y)+\ip{y-\proj_Cy}{x-y}+\tfrac12\norm{x-y}^2 ,
\]
that is, $\varphi(y)\ge\varphi(x)+\ip{y-\proj_Cy}{y-x}-\tfrac12\norm{y-x}^2$. By Theorem
\ref{thm:proj}(3) the map $y\mapsto y-\proj_Cy$ is continuous, so
$\ip{y-\proj_Cy}{y-x}=\ip{g}{y-x}+o(\norm{y-x})$. Combining the two bounds,
$\varphi(y)=\varphi(x)+\ip{g}{y-x}+o(\norm{y-x})$.
\end{proof}

\begin{remark}
Proposition \ref{prop:dist} is the one place this module touches convex functions. It is
included because M3 will recognize $\varphi$ as the \emph{Moreau envelope} of the indicator
function of $C$ and $\proj_C$ as its \emph{proximal operator}, and will prove the gradient
formula for a general convex function in place of an indicator. The proof above is that general
proof, specialized.
\end{remark}

Now the closed forms. Each is verified the same way: guess $p$, then check the inequality of
Theorem \ref{thm:proj}(2). \emph{No optimization is needed once the characterization is in
hand}, and this is the practical payoff of part (2).

\begin{proposition}[Three easy projections]
\label{prop:easyproj}
\begin{enumerate}
\item Halfspace $C=\{y:\ip{a}{y}\le\beta\}$, $a\ne0$:
$\proj_Cx=x-\dfrac{(\ip{a}{x}-\beta)_+}{\norm a^2}\,a$.
\item Affine set $C=\{y:Ay=b\}$ with $A$ of full row rank:
$\proj_Cx=x-A^\top(AA^\top)^{-1}(Ax-b)$.
\item Ball $C=\bar B(c,r)$: $\proj_Cx=x$ if $\norm{x-c}\le r$, and
$c+r\frac{x-c}{\norm{x-c}}$ otherwise. Box $C=\prod_i[l_i,u_i]$: componentwise clipping.
\end{enumerate}
\end{proposition}

\begin{proof}
(1) If $\ip{a}{x}\le\beta$ then $x\in C$ and $p=x$ trivially. Otherwise let
$p=x-\theta a$ with $\theta=(\ip{a}{x}-\beta)/\norm a^2>0$; then $\ip{a}{p}=\beta$ so $p\in C$,
and for $y\in C$, $\ip{x-p}{y-p}=\theta\ip{a}{y-p}=\theta(\ip{a}{y}-\beta)\le0$. Apply Theorem
\ref{thm:proj}(2).
(2) With $p=x-A^\top(AA^\top)^{-1}(Ax-b)$ one checks $Ap=b$, so $p\in C$; and for $y\in C$,
$x-p=A^\top\mu$ with $\mu=(AA^\top)^{-1}(Ax-b)$, so
$\ip{x-p}{y-p}=\ip{\mu}{Ay-Ap}=\ip{\mu}{b-b}=0\le0$. \emph{The inequality holds with equality,
which is exactly the statement that for an affine set the characterization degenerates to
orthogonality} --- the normal equations of least squares.
(3) Analogous; both are left as E3.2.
\end{proof}

The simplex is the one that matters for the rest of the plan, and it is also the one where the
obvious guess is wrong.

\begin{proposition}[Projection onto the probability simplex]
\label{prop:simplex}
Let $x\in\R^n$ and let $\theta$ be the unique solution of $\sum_{i=1}^n(x_i-\theta)_+=1$. Then
$\proj_{\Simplex_n}(x)_i=(x_i-\theta)_+$.
\end{proposition}

\begin{proof}
\emph{$\theta$ exists and is unique.} $g(\theta)=\sum_i(x_i-\theta)_+$ is continuous and
nonincreasing, $g(\theta)=0$ for $\theta\ge\max_ix_i$, and $g(\theta)\to+\infty$ as
$\theta\to-\infty$; moreover $g$ is strictly decreasing wherever it is positive. Hence
$g(\theta)=1$ has exactly one solution.

\emph{The formula.} Put $p_i=(x_i-\theta)_+$. Then $p\ge0$ and $\sum_ip_i=g(\theta)=1$, so
$p\in\Simplex_n$. Observe
\[
x_i-p_i=\begin{cases}\theta,& x_i\ge\theta\ \ (\text{i.e.\ }p_i=x_i-\theta\ge0),\\
x_i<\theta,& x_i<\theta\ \ (\text{i.e.\ }p_i=0),\end{cases}
\qquad\text{so in all cases } x_i-p_i\le\theta,
\]
with equality whenever $p_i>0$. Now take any $y\in\Simplex_n$:
\[
\ip{x-p}{y-p}=\sum_i(x_i-p_i)y_i-\sum_i(x_i-p_i)p_i .
\]
In the second sum only indices with $p_i>0$ contribute, and there $x_i-p_i=\theta$, so that sum
equals $\theta\sum_ip_i=\theta$. In the first sum, $y_i\ge0$ and $x_i-p_i\le\theta$, so it is at
most $\theta\sum_iy_i=\theta$. Hence $\ip{x-p}{y-p}\le\theta-\theta=0$, and Theorem
\ref{thm:proj}(2) identifies $p$ as the projection.
\end{proof}

\begin{remark}
The formula subtracts a \emph{common threshold} $\theta$ and then clips. The map that clips
first and renormalizes, $x\mapsto(x)_+/\sum_i(x_i)_+$, is a different map, and the worked example
below shows how different. Sorting $x$ makes $\theta$ computable in $O(n\log n)$: with
$u_1\ge\dots\ge u_n$ the sorted entries and $\rho=\max\{j:u_j-\frac1j(\sum_{i\le j}u_i-1)>0\}$,
one has $\theta=\frac1\rho(\sum_{i\le\rho}u_i-1)$.
\end{remark}

\subsection*{Worked example}

Let $x=(0.5,\ 0.6,\ -0.2,\ 1.1)\in\R^4$. Sorted descending, $u=(1.1,0.6,0.5,-0.2)$ with partial
sums $1.1,\ 1.7,\ 2.2,\ 2.0$. Test the condition $u_j-\frac1j(\sum_{i\le j}u_i-1)>0$:
\[
\begin{aligned}
j=1:&\ \ 1.1-0.1=1.0>0, &\qquad j=2:&\ \ 0.6-0.35=0.25>0,\\
j=3:&\ \ 0.5-0.4=0.1>0, &\qquad j=4:&\ \ -0.2-0.25=-0.45<0 .
\end{aligned}
\]
So $\rho=3$ and $\theta=\tfrac13(2.2-1)=0.4$. The projection is
\[
\proj_{\Simplex}(x)=\big((0.5-0.4)_+,\ (0.6-0.4)_+,\ (-0.2-0.4)_+,\ (1.1-0.4)_+\big)
=(0.1,\ 0.2,\ 0,\ 0.7),
\]
which is nonnegative and sums to $1$. Check the theorem: $x-p=(0.4,0.4,-0.2,0.4)$, and indeed
every entry is $\le\theta=0.4$, with equality exactly at the coordinates where $p_i>0$.

Now the naive map: $(x)_+=(0.5,0.6,0,1.1)$ with sum $2.2$, giving
$(0.2273,\ 0.2727,\ 0,\ 0.5)$. The two answers differ by
\[
\norm{(0.1,0.2,0,0.7)-(0.2273,0.2727,0,0.5)}
=\norm{(-0.1273,\ -0.0727,\ 0,\ 0.2)}\approx0.248 ,
\]
which is a quarter of the diameter of the simplex --- not a rounding difference. \emph{The
naive map is not the projection, and now you have seen the number.} Note also where they differ
in kind: the projection preserves the \emph{gaps} $x_i-x_j$ among surviving coordinates, while
renormalization preserves their \emph{ratios}.

\subsection*{Where this is used later}

$\proj_C$ is not a computational convenience introduced here for later use. It is the
separating vector of S4, written $x-\proj_Cx$; the stationarity condition of M8, written
$x=\proj_C(x-t\nabla f(x))$; the Moreau decomposition of S5; and the Farkas certificate of S6.
With $C$ an affine subspace, Theorem \ref{thm:proj}(2) becomes the orthogonality $\ip{x-p}{y-p}=0$
of Proposition \ref{prop:easyproj}(2), which is the normal equations, which is the whole of
MMC201's module M2. One theorem, four appearances.

\mpart{Reinforcement learning.} $\proj_\Simplex$ is the operation every projected
policy-gradient step performs, and Proposition \ref{prop:simplex} is why it costs $O(n\log n)$
per state rather than a quadratic program. Read what the formula says: it zeroes every action
whose score falls below a common threshold. \emph{A projected policy step is sparsifying}, and
can reach --- and stop at --- a deterministic policy in finitely many steps, whereas a softmax
parameterization assigns positive probability to every action forever. That contrast between
the Euclidean and the entropic geometry is the subject of M9, and the clip-and-renormalize error
above is the standard way it gets silently broken in code.

\subsection*{Proof of concept: the projection library}

Implement \texttt{project} for the halfspace, the affine set $\{y:Ay=b\}$, the ball, the box,
the second-order cone $K=\{(z,t):\norm z\le t\}$ and $\Simplex_n$, attaching each to the set
objects of S1. For the second-order cone the formula is: $(z,t)$ if $\norm z\le t$; $0$ if
$\norm z\le-t$; and otherwise $\frac{\norm z+t}{2\norm z}(z,\norm z)$ --- derive it in E3.4
rather than trusting it.

\emph{Validation.}
\begin{enumerate}
\item The characterization: for each set, $10^3$ random exterior $x$ and $10^4$ sampled
$y\in C$, $\max\ip{x-\proj_Cx}{y-\proj_Cx}\le10^{-12}$.
\item Firm nonexpansiveness and nonexpansiveness on $10^4$ random pairs: report the maximum
violation of each (both $\le10^{-12}$), and the maximum ratio
$\norm{\proj_Cx-\proj_Cz}/\norm{x-z}$, which must not exceed $1$ and should be found near $1$
for pairs already in $C$.
\item Each closed form agrees with a \texttt{cvxpy} solve of $\min\norm{y-x}^2$ s.t.\ $y\in C$
to $10^{-8}$ --- the solver's tolerance, not machine precision --- on $500$ random instances per
set.
\item Proposition \ref{prop:simplex}: the output lies in $\Simplex_n$ to $10^{-14}$; and
timing over $n=10^2,\dots,10^6$ gives a fitted slope of $\log(\text{time})$ against $\log n$ of
$1+o(1)$, not $2$. \emph{Report the fitted slope.}
\item \textbf{$\proj_C$ is not linear.} Exhibit $C$, $x$, $z$, $\lambda$ with
$\proj_C(\lambda x+(1-\lambda)z)\ne\lambda\proj_Cx+(1-\lambda)\proj_Cz$, and report the gap.
Then verify that for $C$ an affine set the same expression vanishes to $10^{-12}$. \emph{Both
halves are required; the second is what makes the first surprising.}
\item \textbf{Reproduce the worked example.} On the $x$ above, confirm $\theta=0.4$ and
$p=(0.1,0.2,0,0.7)$ exactly, and that the clip-and-renormalize map returns
$(0.2273,0.2727,0,0.5)$ at distance $0.248$ from it. Then report the maximum discrepancy between
the two maps over $10^4$ random $x$ in $\R^{50}$.
\end{enumerate}

\emph{Deliverable.} \texttt{s3.py} and \texttt{project} in \texttt{m1lib.py}. \textbf{Keep it:}
S4 builds separating hyperplanes from it, S5 builds Moreau's decomposition from it, S6's Farkas
certificate is one of its outputs, and M8 and M9 import it unchanged.

\subsection*{Exercises}
\begin{enumerate}[label=E3.\arabic*]
\item Reproduce Theorem \ref{thm:proj}, all three parts, with the document closed. Mark the one
line that uses closedness and the one line that uses convexity.
\item Prove Proposition \ref{prop:easyproj}(3), and verify that in part (2) the characterizing
inequality holds with equality. Then write the normal equations of least squares and identify
them with that equality. \textbf{[This is MMC201's M2.]}
\item Prove Proposition \ref{prop:simplex} without looking, and state precisely where
$y\in\Simplex_n$ (rather than merely $y\ge0$) was used.
\item Derive the second-order cone projection stated above from Theorem \ref{thm:proj}(2).
Hint: by symmetry the answer has the form $\alpha(z,\norm z)$ on the ``otherwise'' branch;
determine $\alpha$ from the characterization.
\item Prove Proposition \ref{prop:dist}. Where does convexity of $C$ enter, and what fails for
$\dist_C$ when $C$ is closed but not convex?
\end{enumerate}

\mpart{Session gate.} Checks 1, 2, 3, 5 and 6 pass, and E3.1 is written out in full.

\clearpage

\section{Session S4. Separation}

\mpart{The question.} A convex set has two descriptions. The \emph{primal} one lists its points:
$C=\conv S$, or $C=\{Ax:x\ge0\}$. The \emph{dual} one lists the constraints its points satisfy:
$C=\{y:\ip{a}{y}\le\alpha\ \text{for every } (a,\alpha)\in\mathcal H\}$. Every duality in this
course is a passage between the two. But why should the second description exist at all, and
how would one produce a constraint that a given exterior point violates?

\mpart{The picture.} Suppose $x\notin C$ with $C$ closed convex. Project: let $p=\proj_Cx$ and
$a=x-p\ne0$. The vector $a$ points from the set towards $x$. By the inequality of S3, no point
of $C$ lies at an angle below $90^\circ$ from $a$ as seen from $p$; so along the direction $a$,
the whole of $C$ sits at or behind the level of $p$, while $x$ sits $\norm a^2$ beyond it. The
hyperplane through the midpoint of the segment $[p,x]$, perpendicular to $a$, separates them ---
and the gap is not merely positive, it is exactly $\norm a^2=\dist_C(x)^2$. \emph{The separating
hyperplane is not an existence statement. It is the output of the projection routine you wrote
in S3.}

That construction requires closedness, and when closedness fails the gap collapses to zero:
the sets can still be separated, but only by a hyperplane that touches both. Keeping those two
cases apart is what the three separation theorems do.

\subsection*{Statements and proofs}

\begin{theorem}[Strong separation of a point]
\label{thm:sep}
Let $C\subseteq\R^n$ be nonempty closed convex and $x\notin C$. Put $a=x-\proj_Cx\ne0$. Then
\[
\ip{a}{c}\ \le\ \ip{a}{x}-\norm a^2\qquad\text{for every }c\in C .
\]
Hence with $\alpha=\ip{a}{x}-\norm a^2$ we have $\sup_{c\in C}\ip{a}{c}\le\alpha<\ip{a}{x}$, and
the gap $\ip{a}{x}-\alpha$ equals $\norm a^2=\dist_C(x)^2$.
\end{theorem}

\begin{proof}
Write $p=\proj_Cx$, so $a=x-p$ and $a\ne0$ because $x\notin C$. Theorem \ref{thm:proj}(2) says
$\ip{a}{c-p}\le0$ for every $c\in C$, i.e. $\ip{a}{c}\le\ip{a}{p}$. Finally
$\ip{a}{p}=\ip{a}{x-a}=\ip{a}{x}-\norm a^2$.
\end{proof}

Two lines. \emph{Write them in a form you can reproduce}: S6 repeats this step with $C$ a cone,
M4 repeats it with $C$ a perturbation epigraph, and M11 repeats it with $C$ the constraint cone
of a linear program. The point of the module is that they repeat \emph{this} step.

\begin{corollary}[The dual description exists]
\label{cor:intersection}
A nonempty closed convex set is the intersection of all the closed halfspaces containing it.
\end{corollary}

\begin{proof}
Let $D$ be that intersection. $C\subseteq D$ by definition. If $x\notin C$, Theorem
\ref{thm:sep} gives a halfspace $\{y:\ip{a}{y}\le\alpha\}$ containing $C$ but not $x$, so
$x\notin D$. Hence $D\subseteq C$.
\end{proof}

\begin{theorem}[Supporting hyperplane]
\label{thm:support}
Let $C$ be convex and $x_0\in\bd(\cl C)$. Then there is $a\ne0$ with
$\ip{a}{c}\le\ip{a}{x_0}$ for every $c\in C$.
\end{theorem}

\begin{proof}
Since $x_0$ is a boundary point of $\cl C$, there are $x_k\notin\cl C$ with $x_k\to x_0$. Put
$a_k=(x_k-\proj_{\cl C}x_k)/\norm{x_k-\proj_{\cl C}x_k}$, a unit vector; by Theorem
\ref{thm:sep} applied to $\cl C$, $\ip{a_k}{c}\le\ip{a_k}{x_k}$ for all $c\in\cl C$. The unit
sphere of $\R^n$ is compact --- \emph{this is the second and last use of finite dimension in
the module} --- so a subsequence has $a_k\to a$ with $\norm a=1$. Passing to the limit in the
inequality gives $\ip{a}{c}\le\ip{a}{x_0}$ for all $c\in\cl C\supseteq C$.
\end{proof}

\begin{theorem}[Strong separation of two sets]
\label{thm:septwo}
Let $C$ be compact convex, $D$ closed convex, both nonempty, with $C\cap D=\emptyset$. Then
there are $a\ne0$ and $\gamma>0$ with
$\sup_{c\in C}\ip{a}{c}\le\inf_{d\in D}\ip{a}{d}-\gamma$.
\end{theorem}

\begin{proof}
Let $E=C-D=\{c-d:c\in C,\ d\in D\}$. It is convex, and it is closed by Proposition
\ref{prop:recession}, whose hypothesis holds because $C$ is bounded so $0^+C=\{0\}$. Since
$C\cap D=\emptyset$ we have $0\notin E$. Apply Theorem \ref{thm:sep} to $E$ and the point $0$:
with $b=0-\proj_E0=-\proj_E0\ne0$,
\[
\ip{b}{e}\le\ip{b}{0}-\norm b^2=-\norm b^2\qquad\text{for all }e\in E .
\]
Writing $e=c-d$ and setting $a=b$, $\gamma=\norm b^2>0$ gives
$\ip{a}{c}\le\ip{a}{d}-\gamma$ for all $c\in C$, $d\in D$, which is the claim.
\end{proof}

The compactness in Theorem \ref{thm:septwo} is not decoration. Here is what is lost without it.

\begin{example}[Two disjoint closed convex sets that cannot be strongly separated]
\label{ex:nosep}
Take $C$ and $D$ from Example \ref{ex:notclosed}: $C=\{(u,v):u>0,v\ge1/u\}$ and
$D=\{0\}\times\R$. Both are closed and convex, and they are disjoint since every point of $C$
has $u>0$. By the computation in Example \ref{ex:notclosed}, $C-D$ is the open right half-plane.
Suppose $a=(a_1,a_2)$ separates: then $\sup_{e\in C-D}\ip{a}{e}<\infty$. Since the second
coordinate of $e$ ranges over all of $\R$, this forces $a_2=0$; since the first ranges over
$(0,\infty)$, it forces $a_1\le0$. Then $\sup_{e\in C-D}\ip{a}{e}=0=\ip{a}{0}$: the supremum is
attained in the limit and the gap is exactly $0$. So $C$ and $D$ can be separated
\emph{properly} --- there is a hyperplane with $C$ weakly on one side and $D$ weakly on the
other --- but never \emph{strongly}, for any $a$ whatever.
\end{example}

For completeness, the theorem that says exactly when proper separation is available; it is the
only statement in this document used without proof besides Proposition \ref{prop:ricl}(2).

\begin{theorem}[Proper separation; stated without proof]
\label{thm:proper}
Two nonempty convex sets $C,D$ can be properly separated if and only if
$\ri C\cap\ri D=\emptyset$.
\end{theorem}

This is why Slater's condition in M4 is a statement about relative interiors, and why S2 had to
come before S4.

\subsection*{Worked example}

Let $C=\bar B(0,1)\subseteq\R^2$ and $x=(3,0)$. Then $\proj_Cx=(1,0)$, so $a=x-\proj_Cx=(2,0)$
and $\norm a^2=4$. Theorem \ref{thm:sep} gives $\alpha=\ip{a}{x}-\norm a^2=6-4=2$, so the
separating hyperplane is $\{y:2y_1=2\}$, i.e.\ $y_1=1$: it is the tangent line to the ball at
the projection. The predicted gap is $\ip{a}{x}-\alpha=4=\dist_C(x)^2=(3-1)^2$. \emph{Check
this equality in the code, not merely that some separation occurred}: many wrong implementations
separate correctly and get the margin wrong, and the margin is the part that carries the
information.

Now the same construction for the polytope $C=\conv\{(0,0),(1,0),(0,1)\}$ and $x=(1,1)$. The
nearest point of $C$ is on the edge $u+v=1$: $\proj_Cx=(0.5,0.5)$, so $a=(0.5,0.5)$,
$\norm a^2=0.5$, and $\alpha=\ip{a}{x}-0.5=1-0.5=0.5$. The hyperplane is
$0.5(y_1+y_2)=0.5$, i.e.\ $y_1+y_2=1$ --- the facet of the triangle that $x$ violates. Note what
happened: \emph{the separation construction recovered a facet of the polytope}, which is how the
same idea will produce cutting planes in M14.

\subsection*{Where this is used later}

Theorem \ref{thm:sep} is used in S5 to prove the bipolar theorem and in S6 to prove Farkas'
lemma; those two uses are the whole remainder of this module. Beyond it: M4 proves strong
duality by separating $(0,p^\star)$ from the perturbation epigraph --- literally this theorem;
M11 obtains LP duality the same way; M14 generates Benders and Gomory cuts by separating the
current point from the convex hull of the feasible set. Example \ref{ex:nosep} is the mechanism
of a nonzero duality gap and of nothing else.

\subsection*{Proof of concept: \texttt{separate(C, x)}}

Return $(a,\alpha)$ with $a=x-\proj_Cx$ and $\alpha=\ip{a}{x}-\norm a^2$, built directly on S3's
\texttt{project}. It is ten lines; the session's work is the mathematics.

\emph{Validation.}
\begin{enumerate}
\item \textbf{The margin is exactly the squared distance.} For each set of S3 and $10^3$
exterior $x$: $\ip{a}{x}-\alpha=\norm{x-\proj_Cx}^2$ to $10^{-12}$. This is the constructive
content of Theorem \ref{thm:sep}, and it is a far stronger check than confirming that some
separation occurred.
\item $\sup\ip{a}{c}\le\alpha+10^{-12}$ over $10^4$ sampled $c\in C$, for every set and every
$x$.
\item Reproduce both worked examples exactly: the ball case must return $a=(2,0)$, $\alpha=2$,
margin $4$; the triangle case $a=(0.5,0.5)$, $\alpha=0.5$, margin $0.5$, and the returned
hyperplane must coincide with the violated facet $y_1+y_2=1$ up to positive scaling.
\item Corollary \ref{cor:intersection}, quantitatively: for $C$ a random polygon in $\R^2$,
generate the halfspaces produced by $m$ exterior points sampled on a large circle, intersect
them, and estimate the area of the intersection by Monte Carlo. Report the areas at
$m=10,10^2,10^3,10^4$ and confirm they decrease monotonically towards the area of $C$ from
above.
\item \textbf{The compactness hypothesis of Theorem \ref{thm:septwo}, failing.} For $C$ and $D$
of Example \ref{ex:nosep}, maximize the margin
$\inf_{d\in D}\ip{a}{d}-\sup_{c\in C}\ip{a}{c}$ over unit $a$ sampled on the circle. Confirm the
best margin found is $0$ to the sampling resolution and never positive, so no strong separation
exists although the sets are disjoint, closed and convex. \textbf{Keep this: a nonzero duality
gap in M4 is this failure and no other.}
\end{enumerate}

\emph{Deliverable.} \texttt{s4.py} and \texttt{separate} in \texttt{m1lib.py}.

\subsection*{Exercises}
\begin{enumerate}[label=E4.\arabic*]
\item Reproduce Theorem \ref{thm:sep} with the document closed, marking where closedness and
where convexity are used --- both are inherited from Theorem \ref{thm:proj}, and you should be
able to say which line of \emph{that} proof each comes from.
\item Prove Corollary \ref{cor:intersection}, and explain with one picture and one sentence why
it fails for a nonconvex closed set. This corollary is the reason a convex set has a dual
description at all, and every duality in the plan is a dual description of something.
\item Prove Theorem \ref{thm:support}, and say precisely why the normalized separating vectors
have a convergent subsequence.
\item Verify Example \ref{ex:nosep} in full and name the failing hypothesis of Theorem
\ref{thm:septwo}. \textbf{[Reused in M4: this is the duality gap.]}
\item Build the table you will need in S7: three rows --- proper, strong, supporting --- and
three columns --- hypothesis, conclusion, counterexample when the hypothesis is dropped. Half a
page.
\end{enumerate}

\mpart{Session gate.} Checks 1, 2, 3 and 5 pass; E4.1 and E4.4 are written out.

\clearpage

\section{Session S5. Cones, polars and Moreau's decomposition}

\mpart{The question.} Separation gives an inequality $\ip{a}{c}\le\alpha$ with some number
$\alpha$ on the right. For a general convex set that number carries information and must be
computed. But a cone is scale-invariant: if $c$ is in it, so is $10^6c$. What does a bounded
linear functional on a cone look like, and what does separation become when the number on the
right can only be $0$?

\mpart{The picture.} If $\ip{a}{c}>0$ for a single point $c$ of a cone $K$, then
$\ip{a}{tc}\to+\infty$, so the supremum of $\ip{a}{\cdot}$ over $K$ is either $0$ or $+\infty$
--- never anything in between. The vectors for which it is $0$ form a set, the polar cone
$K^\circ$, and separation from a cone always produces a member of it. \emph{This is why cones
give exact dualities where general convex sets give only inequalities}, and it is the reason
linear programming has no duality gap while general convex programming needs a hypothesis.

The second picture is a decomposition. For a subspace $L$, every vector splits uniquely as
something in $L$ plus something in $L^\perp$, and the two pieces are the two projections. Replace
the subspace by a closed convex cone and the orthogonal complement by the polar cone: the same
statement is true, with orthogonality of the two pieces surviving intact. That is Moreau's
theorem, and it costs three lines given S3.

\subsection*{Statements and proofs}

Recall Definition \ref{def:polar}: $S^\circ=\{y:\ip{y}{s}\le0\ \forall s\in S\}$. Note that
$S^\circ$ is always a closed convex cone, being an intersection of halfspaces through the
origin, whatever $S$ is.

\begin{theorem}[Bipolar]
\label{thm:bipolar}
For any nonempty $S\subseteq\R^n$, $S^{\circ\circ}=\cl\conv\cone(S)$. In particular a nonempty
set $K$ satisfies $K^{\circ\circ}=K$ if and only if $K$ is a closed convex cone.
\end{theorem}

\begin{proof}
Write $K=\cl\conv\cone(S)$; it is a nonempty closed convex cone. First, $S^\circ=K^\circ$:
the inclusion $K^\circ\subseteq S^\circ$ holds since $S\subseteq K$; conversely if
$\ip{y}{s}\le0$ for all $s\in S$ then $\ip{y}{\cdot}\le0$ on all nonnegative combinations of
elements of $S$ by linearity, and on their closure by continuity, so $y\in K^\circ$.

$(\supseteq)$ Let $x\in K$ and $y\in S^\circ=K^\circ$. Then $\ip{y}{x}\le0$, and since $y$ was
arbitrary in $S^\circ$, $x\in S^{\circ\circ}$.

$(\subseteq)$ Let $x\notin K$. Since $K$ is nonempty closed convex, Theorem \ref{thm:sep} gives
$a\ne0$ with $\ip{a}{k}\le\ip{a}{x}-\norm a^2$ for all $k\in K$; in particular
$\sup_{k\in K}\ip{a}{k}<\infty$. If $\ip{a}{k_0}>0$ for some $k_0\in K$, then $tk_0\in K$ for
all $t>0$ and $\ip{a}{tk_0}\to+\infty$, a contradiction. Hence $\ip{a}{k}\le0$ for all $k\in K$,
that is, $a\in K^\circ=S^\circ$. Also $0\in K$, so $0\le\ip{a}{x}-\norm a^2$, giving
$\ip{a}{x}\ge\norm a^2>0$. Thus $a\in S^\circ$ with $\ip{a}{x}>0$, so $x\notin S^{\circ\circ}$.

For the last sentence: if $K^{\circ\circ}=K$ then $K$ is a closed convex cone because every
polar is; conversely if $K$ is a nonempty closed convex cone then $\cl\conv\cone(K)=K$.
\end{proof}

\begin{theorem}[Moreau decomposition]
\label{thm:moreau}
Let $K\subseteq\R^n$ be a closed convex cone. Every $x\in\R^n$ has a unique decomposition
\[
x=u+v,\qquad u\in K,\quad v\in K^\circ,\quad \ip{u}{v}=0 ,
\]
and it is given by $u=\proj_Kx$, $v=\proj_{K^\circ}x$.
\end{theorem}

\begin{proof}
\emph{Construction.} Let $p=\proj_Kx$ and $w=x-p$. By Theorem \ref{thm:proj}(2),
$\ip{w}{y-p}\le0$ for all $y\in K$. Taking $y=0\in K$ gives $\ip{w}{p}\ge0$; taking
$y=2p\in K$ (a cone) gives $\ip{w}{p}\le0$. Hence $\ip{w}{p}=0$. Then for any $y\in K$,
\[
\ip{w}{y}=\ip{w}{y-p}+\ip{w}{p}\le0 ,
\]
so $w\in K^\circ$. Thus $x=p+w$ with $p\in K$, $w\in K^\circ$, $\ip{p}{w}=0$.

\emph{$w$ is the projection onto $K^\circ$.} For any $z\in K^\circ$,
\[
\ip{x-w}{z-w}=\ip{p}{z-w}=\ip{p}{z}-\ip{p}{w}=\ip{p}{z}\le0 ,
\]
the last step because $p\in K$ and $z\in K^\circ$. By Theorem \ref{thm:proj}(2),
$w=\proj_{K^\circ}x$.

\emph{Uniqueness.} Suppose $x=u+v$ with $u\in K$, $v\in K^\circ$, $\ip{u}{v}=0$. For $y\in K$,
\[
\ip{x-u}{y-u}=\ip{v}{y}-\ip{v}{u}=\ip{v}{y}\le0 ,
\]
so $u=\proj_Kx$ by Theorem \ref{thm:proj}(2), and then $v=x-u$ is determined.
\end{proof}

\begin{corollary}
\label{cor:subspace}
If $K=L$ is a subspace then $L^\circ=L^\perp$ and Theorem \ref{thm:moreau} is the orthogonal
decomposition $x=\proj_Lx+\proj_{L^\perp}x$.
\end{corollary}

\begin{proof}
For a subspace, $l\in L$ implies $-l\in L$, so $\ip{y}{l}\le0$ for all $l\in L$ forces
$\ip{y}{l}=0$; hence $L^\circ=L^\perp$.
\end{proof}

\begin{proposition}[Polars of the cones we need]
\label{prop:polarex}
\begin{enumerate}
\item $(\R^n_+)^\circ=\R^n_-$.
\item $\cone(A)^\circ=\{y:A^\top y\le0\}$ for any $A\in\R^{m\times n}$.
\item The second-order cone $K=\{(z,t)\in\R^{n-1}\times\R:\norm z\le t\}$ satisfies
$K^\circ=-K$; equivalently $K^*=K$, so $K$ is self-dual.
\end{enumerate}
\end{proposition}

\begin{proof}
(1) If $y\le0$ then $\ip{y}{x}\le0$ for $x\ge0$. Conversely taking $x=e_i$ gives $y_i\le0$.

(2) $y\in\cone(A)^\circ$ iff $\ip{y}{Ax}\le0$ for all $x\ge0$, iff $\ip{A^\top y}{x}\le0$ for
all $x\ge0$, iff $A^\top y\le0$ by part (1).

(3) Let $(w,s)\in K^*$, i.e.\ $\ip{w}{z}+st\ge0$ for all $(z,t)\in K$. Taking $t=1$ and
$z=-w/\norm w$ (legitimate since $\norm z=1\le t$) gives $-\norm w+s\ge0$, i.e.\ $s\ge\norm w$,
so $(w,s)\in K$. Conversely if $\norm w\le s$ and $\norm z\le t$ then by Cauchy--Schwarz
$\ip{w}{z}+st\ge-\norm w\norm z+st\ge-st+st=0$. Hence $K^*=K$.
\end{proof}

\begin{remark}
The companion identity $\{x:Ax\le0\}^\circ=\cone(A^\top)$ is \emph{not} proved here. By part (2)
applied to $A^\top$ we have $\{x:Ax\le0\}=\cone(A^\top)^\circ$, so taking polars and using
Theorem \ref{thm:bipolar} gives $\{x:Ax\le0\}^\circ=\cl\cone(A^\top)$ --- and the closure can be
removed only once one knows that $\cone(A^\top)$ is closed. That is Lemma \ref{lem:closedcone},
and the resulting identity is Farkas' lemma. \emph{So the missing half of this proposition is
exactly the agenda of S6.}
\end{remark}

\begin{definition}[Normal cone]
\label{def:normal}
For $C$ convex and $x\in C$, $\Nc_C(x)=\{v:\ip{v}{y-x}\le0\ \forall y\in C\}$. By Theorem
\ref{thm:proj}(2), $v\in\Nc_C(x)$ if and only if $x=\proj_C(x+v)$.
\end{definition}

\begin{proposition}[The normal cone of the simplex]
\label{prop:normalsimplex}
For $x\in\Simplex_n$, $\Nc_{\Simplex_n}(x)=\{v:\ \{i:x_i>0\}\subseteq\argmax_jv_j\}$.
\end{proposition}

\begin{proof}
$v\in\Nc(x)$ iff $\ip{v}{y}\le\ip{v}{x}$ for all $y\in\Simplex_n$, iff
$\max_{y\in\Simplex_n}\ip{v}{y}=\ip{v}{x}$. The maximum of a linear functional over the simplex
is $\max_jv_j$, attained exactly at the distributions supported on $\argmax_jv_j$. So the
condition is that $x$ is supported there.
\end{proof}

\subsection*{Worked example}

Take $K=\R^2_+$ and $x=(3,-2)$. Then $\proj_Kx=(3,0)$ by componentwise clipping, and
$K^\circ=\R^2_-$ with $\proj_{K^\circ}x=(0,-2)$. Indeed $(3,0)+(0,-2)=(3,-2)=x$ and
$\ip{(3,0)}{(0,-2)}=0$: Theorem \ref{thm:moreau} verified by hand.

Now the normal cone of $\Simplex_3$, from Proposition \ref{prop:normalsimplex}. At the interior
point $x=(\tfrac13,\tfrac13,\tfrac13)$ the support is everything, so $\argmax_jv_j$ must be all
of $\{1,2,3\}$: $\Nc(x)=\{v:v_1=v_2=v_3\}=\R\cdot\one$, a line. At the vertex $x=e_1$ the
support is $\{1\}$, so the condition is $v_1\ge v_2$ and $v_1\ge v_3$: $\Nc(e_1)$ is a
\emph{full-dimensional} cone in $\R^3$. At the edge point $x=(\tfrac12,\tfrac12,0)$ the
condition is $v_1=v_2\ge v_3$, a two-dimensional cone. The normal cone grows as the face
shrinks, which is the general pattern and the reason faces and duality are the same subject.

\subsection*{Where this is used later}

Theorem \ref{thm:bipolar} and Proposition \ref{prop:polarex}(2) are the two ingredients S6
combines with Lemma \ref{lem:closedcone} to get Farkas. Definition \ref{def:normal} is how M4
states optimality --- $-\nabla f(x^\star)\in\Nc_C(x^\star)$ --- and how M8 states stationarity
of the projected gradient method. Theorem \ref{thm:moreau} is the conic case of the general
decomposition $x=\prox_f(x)+\prox_{f^*}(x)$ that M3 proves for an arbitrary closed convex
function; meeting it here costs three lines and makes M3's version recognizable.

\mpart{Reinforcement learning.} Proposition \ref{prop:normalsimplex} is why stationarity of a
projected policy-gradient method is a set inclusion rather than an equation. At a deterministic
policy --- a vertex of the simplex --- the normal cone is full-dimensional, so an entire cone of
gradients is consistent with stationarity: the method can \emph{stop} at a deterministic policy,
and generically does. A softmax parameterization has no boundary and therefore no normal cones,
and approaches determinism only as its parameters diverge. The same contrast in a different
guise --- Euclidean versus entropic geometry --- is the subject of M9.

\subsection*{Proof of concept: polars and the conic decomposition}

Implement \texttt{polar} for a polyhedral cone in either representation, using Proposition
\ref{prop:polarex}(2), and \texttt{moreau\_decompose(K, x)} using S3's \texttt{project}.

\emph{Validation.}
\begin{enumerate}
\item Theorem \ref{thm:moreau} for $K=\R^n_+$, the second-order cone, a random polyhedral cone,
and a random subspace, over $10^4$ random $x$:
$\norm{x-\proj_Kx-\proj_{K^\circ}x}\le10^{-10}$ and
$|\ip{\proj_Kx}{\proj_{K^\circ}x}|\le10^{-12}$. For the subspace, additionally
$\proj_K+\proj_{K^\perp}=I$ to $10^{-14}$ and both maps are linear, checked as in S3 check 5 ---
which is Corollary \ref{cor:subspace}.
\item Proposition \ref{prop:polarex}(3): over $10^5$ random pairs $y,c$ in the second-order cone,
$\ip{y}{c}\ge-10^{-12}$; and the same for $\R^n_+$.
\item Theorem \ref{thm:bipolar}, computationally: take a random cone by generators, compute its
polar as an inequality description, then the bipolar back as generators, and confirm the two
generator sets span the same cone by testing membership in both directions with
\texttt{linprog} to $10^{-9}$. \emph{Include one cone that is not full-dimensional and one that
is not pointed}, since those are the cases where a careless implementation silently drops a
lineality direction.
\item Reproduce the worked example: $\Nc_{\Simplex_3}$ at an interior point, a vertex and an
edge point, verified by sampling $v$ and checking $\ip{v}{y-x}\le10^{-12}$ over $y\in\Simplex_3$.
Report the dimension of each normal cone, estimated from the rank of a sample of its members;
the answers must be $1$, $3$ and $2$.
\end{enumerate}

\emph{Deliverable.} \texttt{s5.py}, with \texttt{polar} and \texttt{moreau\_decompose} in
\texttt{m1lib.py}.

\subsection*{Exercises}
\begin{enumerate}[label=E5.\arabic*]
\item Reproduce Theorem \ref{thm:bipolar} with the document closed, isolating the sentence in
which the cone property --- ``the supremum is $0$ or $+\infty$'' --- is used.
\item Reproduce Theorem \ref{thm:moreau} and deduce Corollary \ref{cor:subspace}.
\textbf{[This is the orthogonal decomposition of MMC201's M2.]}
\item Prove Proposition \ref{prop:polarex} in full, from Definition \ref{def:polar}, with no
appeal to a table.
\item Write out the Remark after Proposition \ref{prop:polarex} as a short argument, ending in
the sentence: ``Farkas' lemma is Theorem \ref{thm:bipolar} applied to $\cone(A)$, and the only
missing ingredient is that $\cone(A)$ is closed.'' \textbf{[This is the agenda of S6.]}
\item Prove Proposition \ref{prop:normalsimplex} and compute $\Nc_{\Simplex_n}(x)$ at a general
point of a proper face. \textbf{[Reused in M8.]}
\end{enumerate}

\mpart{Session gate.} Checks 1, 3 and 4 pass; E5.2 and E5.4 are written out.

\clearpage

\section{Session S6. Closed cones and Farkas' lemma}

\mpart{The question.} You are handed $A$ and $b$ and told that $Ax=b,\ x\ge0$ has no solution.
How would you \emph{prove} it to someone who does not trust you? Exhibiting your failure to find
$x$ proves nothing. What is wanted is a short certificate: an object that can be checked in a
few arithmetic operations and that makes the nonexistence evident. Does one always exist?

\mpart{The picture.} ``$Ax=b$ for some $x\ge0$'' says exactly $b\in\cone(A)$. So the question is
geometric: $b$ is outside a cone, and we want a certificate of outsideness. S4 supplies one ---
a separating hyperplane --- and S5 says that for a cone the hyperplane passes through the origin,
so the certificate is a single vector $y$ with $\ip{y}{c}\le0$ for all $c\in\cone(A)$, i.e.\
$A^\top y\le0$, and $\ip{y}{b}>0$. Flip signs and that is the second alternative.

But there is a gap in that argument, and it is the whole content of this session. Theorem
\ref{thm:sep} requires $\cone(A)$ to be \emph{closed}, and by S2 we know that images of closed
sets under linear maps are not closed in general --- $\cone(A)$ is precisely the image of
$\R^n_+$ under $A$. So the geometry gives the certificate only after one proves that this
particular image is closed. That is Lemma \ref{lem:closedcone}, it is true because there are
only \emph{finitely many} generators, and it is the reason Farkas' lemma is a statement that
exactly one of two systems holds rather than merely that they cannot both hold.

\subsection*{Statements and proofs}

\begin{lemma}[Finitely generated cones are closed]
\label{lem:closedcone}
For any $A\in\R^{m\times n}$, the cone $\cone(A)=\{Ax:x\ge0\}$ is closed.
\end{lemma}

\begin{proof}
For a subset $I\subseteq\{1,\dots,n\}$ of column indices, write $A_I$ for the submatrix of those
columns, and let $\mathcal I$ be the collection of $I$ for which the columns of $A_I$ are
linearly independent. By Corollary \ref{cor:conecara}, every element of $\cone(A)$ lies in
$\cone(A_I)$ for some $I\in\mathcal I$, so
\[
\cone(A)=\bigcup_{I\in\mathcal I}\cone(A_I) .
\]
Fix $I\in\mathcal I$ and let $A_Ix_k\to w$ with $x_k\ge0$. Since the columns of $A_I$ are
independent, $A_I$ is injective, so $A_I^\top A_I$ is invertible and
$x_k=(A_I^\top A_I)^{-1}A_I^\top(A_Ix_k)\to(A_I^\top A_I)^{-1}A_I^\top w=:x_\infty$ by
continuity. Then $x_\infty\ge0$ as a limit of nonnegative vectors, and $A_Ix_\infty=w$ by
continuity; so $w\in\cone(A_I)$ and $\cone(A_I)$ is closed. Finally $\mathcal I$ is
\emph{finite}, and a finite union of closed sets is closed.
\end{proof}

\begin{remark}
Two finiteness facts were used: finitely many index subsets, and injectivity of $A_I$ obtained
from linear independence. Both fail for a cone generated by infinitely many vectors, and the
conclusion fails with them --- S7 exhibits a convex cone that is not closed. In M4 the cone in
question is a tangent cone, which need not be polyhedral; \emph{a constraint qualification is
precisely a hypothesis restoring the conclusion of this lemma}, and that is the single most
useful thing to carry from S6 into M4.
\end{remark}

\begin{theorem}[Farkas]
\label{thm:farkas}
Let $A\in\R^{m\times n}$ and $b\in\R^m$. Exactly one of the following systems has a solution:
\begin{enumerate}
\item[(I)] $x\in\R^n$ with $Ax=b$ and $x\ge0$;
\item[(II)] $y\in\R^m$ with $A^\top y\ge0$ and $\ip{b}{y}<0$.
\end{enumerate}
\end{theorem}

\begin{proof}
\emph{Not both.} Suppose $x$ solves (I) and $y$ solves (II). Then
\[
\ip{b}{y}=\ip{Ax}{y}=\ip{x}{A^\top y}\ge0 ,
\]
since $x\ge0$ and $A^\top y\ge0$, contradicting $\ip{b}{y}<0$. \emph{This is one line, it uses
nothing but the definitions, and it is weak duality in its simplest form.}

\emph{At least one.} Suppose (I) has no solution, i.e.\ $b\notin\cone(A)$. By Lemma
\ref{lem:closedcone} the set $\cone(A)$ is closed; it is convex and nonempty; so Theorem
\ref{thm:sep} applies and yields $a=b-\proj_{\cone(A)}b\ne0$ with
\[
\ip{a}{c}\ \le\ \ip{a}{b}-\norm a^2\qquad\text{for every }c\in\cone(A) .
\]
In particular $\sup_{c\in\cone(A)}\ip{a}{c}<\infty$, and since $\cone(A)$ is a cone this
supremum is $0$ (as in the proof of Theorem \ref{thm:bipolar}: a single $c$ with
$\ip{a}{c}>0$ would make it $+\infty$). Taking $c=0$ in the display gives
$0\le\ip{a}{b}-\norm a^2$, so $\ip{a}{b}\ge\norm a^2>0$. And $\ip{a}{Ax}\le0$ for all $x\ge0$
means $A^\top a\le0$. Now put $y=-a$: then $A^\top y\ge0$ and
$\ip{b}{y}=-\ip{a}{b}\le-\norm a^2<0$, so $y$ solves (II).
\end{proof}

\mpart{The one line that matters.} In the second half of that proof, everything is a definition
chase except the appeal to Theorem \ref{thm:sep}, and that appeal is legitimate \emph{only
because of Lemma \ref{lem:closedcone}}. Mark it. The entire difference between the two halves of
the proof is that the first needs nothing and the second needs the set to be closed --- which is
the same as saying that weak duality is free and strong duality is not.

\begin{corollary}[Inequality form]
\label{cor:farkasineq}
Exactly one of the following has a solution: $x$ with $Ax\le b$; or $y\ge0$ with $A^\top y=0$
and $\ip{b}{y}<0$.
\end{corollary}

\begin{proof}
Write $x=x^+-x^-$ with $x^\pm\ge0$ and introduce a slack $s\ge0$: the system $Ax\le b$ is
solvable iff
\[
\begin{bmatrix}A&-A&I\end{bmatrix}
\begin{bmatrix}x^+\\x^-\\s\end{bmatrix}=b,\qquad
(x^+,x^-,s)\ge0
\]
is solvable. Apply Theorem \ref{thm:farkas} to $\tilde A=[A\ \ {-A}\ \ I]$: the alternative is
$y$ with $\tilde A^\top y\ge0$ and $\ip{b}{y}<0$, and $\tilde A^\top y\ge0$ reads
$A^\top y\ge0$, $-A^\top y\ge0$, $y\ge0$, i.e.\ $A^\top y=0$ and $y\ge0$.
\end{proof}

\begin{theorem}[Gordan]
\label{thm:gordan}
Exactly one of the following has a solution: $x$ with $Ax<0$ (strictly, componentwise); or
$y\ge0$, $y\ne0$, with $A^\top y=0$.
\end{theorem}

\begin{proof}
First, $Ax<0$ has a solution iff $Ax\le-\one$ has one: if $Ax<0$, put
$\sigma=\min_i(-(Ax)_i)>0$ and use $x/\sigma$; conversely $Ax\le-\one$ gives $Ax<0$. Apply
Corollary \ref{cor:farkasineq} with right-hand side $-\one$: exactly one of $Ax\le-\one$, and
$y\ge0$ with $A^\top y=0$, $\ip{-\one}{y}<0$. The last condition is $\sum_iy_i>0$, which for
$y\ge0$ is the same as $y\ne0$.
\end{proof}

Stiemke's and Motzkin's theorems have the same shape and are obtained the same way; the pattern
worth carrying forward is that \emph{each theorem of the alternative is Theorem
\ref{thm:farkas} for a different cone, and the cone can be read off the statement.} That is
worth more than memorizing four statements.

\subsection*{Worked example}

Let
\[
A=\begin{bmatrix}1&1\\1&1\end{bmatrix},\qquad b=\begin{bmatrix}1\\2\end{bmatrix} .
\]
Here $\cone(A)=\{s(1,1):s\ge0\}$, a ray, and $b$ is not on it, so (I) has no solution. Run the
proof and extract the certificate.

The projection of $b$ onto the ray in direction $\hat u=(1,1)/\sqrt2$ is
$\proj b=\ip{b}{\hat u}\hat u=\tfrac{3}{\sqrt2}\cdot\tfrac{1}{\sqrt2}(1,1)=(1.5,1.5)$ --- the
inner product $\tfrac3{\sqrt2}$ is positive, so the projection lands in the ray rather than at
the apex. Hence
\[
a=b-\proj b=(-0.5,\ 0.5),\qquad y=-a=(0.5,\ -0.5) .
\]
Check the certificate: $A^\top y=(0.5-0.5,\ 0.5-0.5)=(0,0)\ge0$ \checkmark, and
$\ip{b}{y}=0.5-1=-0.5<0$ \checkmark. So $y$ proves that no $x\ge0$ has $Ax=b$.

Check the margin: $-\ip{b}{y}=0.5$ and $\norm{b-\proj b}^2=0.25+0.25=0.5$. \emph{They agree,
and they must}, because this is Theorem \ref{thm:sep}'s margin identity carried through Theorem
\ref{thm:farkas}. Verifying that equality in code is the single most informative check in this
module: it tests the proof rather than the statement.

\subsection*{Where this is used later}

Theorem \ref{thm:farkas} is quoted three more times in the plan, each time with $\cone(A)$
replaced by another closed convex cone:
\begin{itemize}
\item in \textbf{M4}, with the linearized cone of the active constraints, giving the
Karush--Kuhn--Tucker conditions --- and, as the Remark after Lemma \ref{lem:closedcone} says, a
constraint qualification is exactly the hypothesis that lets the argument run;
\item in \textbf{M11}, with the constraint cone of a linear program, giving strong duality with
no qualification needed, because Lemma \ref{lem:closedcone} supplies closedness for free in the
polyhedral case. \emph{That sentence is the reason linear programming duality is cleaner than
convex duality};
\item in \textbf{M14}, with the dual feasible cone of a Benders subproblem, where the extreme
rays produced by the alternative are the feasibility cuts.
\end{itemize}

\mpart{Reinforcement learning.} In M11 the Bellman inequalities $V\ge r_a+\gamma P_aV$ define
the feasible set of the linear programming formulation of a Markov decision process. Theorem
\ref{thm:farkas} is what certifies that a modified system --- for instance one with an added
constraint encoding a safety budget --- is infeasible, and the certificate $y$ is, up to
normalization, an occupancy measure. Infeasibility certificates and occupancy measures are the
same objects seen from two sides.

\subsection*{Proof of concept: the Farkas certificate}

Implement \texttt{farkas\_certificate(A, b)} returning either $x\ge0$ with $Ax=b$, or $y$ with
$A^\top y\ge0$ and $\ip{b}{y}<0$ --- never both and never neither.

\emph{Validation.}
\begin{enumerate}
\item On $10^3$ random feasible instances (draw $x_0\ge0$, set $b=Ax_0$): the routine returns
$x\ge-10^{-12}$ with $\norm{Ax-b}\le10^{-9}$, and reports no certificate.
\item On $10^3$ random infeasible instances (draw $y_0$, choose $A$ with $A^\top y_0\ge0$ and
$b$ with $\ip{b}{y_0}<0$): the returned $y$ satisfies $A^\top y\ge-10^{-10}$ componentwise and
$\ip{b}{y}\le-10^{-6}$, and independently \texttt{scipy.optimize.linprog} reports the primal
infeasible. \emph{Both directions must be checked on the same instance;} a routine that only
ever confirms the branch it took has not been tested.
\item \textbf{The certificate is the projection.} Verify that
$y=-(b-\proj_{\cone(A)}b)$, computed with the S3 and S5 code, is a valid certificate, and that
its margin $-\ip{b}{y}$ equals $\norm{b-\proj_{\cone(A)}b}^2$ to $10^{-9}$. \emph{This is the
proof of Theorem \ref{thm:farkas} executed rather than read.}
\item Reproduce the worked example exactly: $\proj b=(1.5,1.5)$, $y=(0.5,-0.5)$,
$A^\top y=(0,0)$, $\ip{b}{y}=-0.5$, margin $0.5$.
\item Theorem \ref{thm:gordan}: over $10^3$ random $A$, exactly one of the two systems is
solvable, with zero disagreements between the two independent solves.
\end{enumerate}

\emph{Deliverable.} \texttt{s6.py} and \texttt{farkas\_certificate} in \texttt{m1lib.py}.
\textbf{Keep it:} M11 asks for the same routine as the infeasibility certificate of a linear
program, and M12's Phase I produces one as a by-product of the simplex method --- at which point
you should check that the two certificates agree up to positive scaling.

\subsection*{Exercises}
\begin{enumerate}[label=E6.\arabic*]
\item Reproduce Lemma \ref{lem:closedcone} with the document closed. Which two finiteness facts
are used? \textbf{[Reused in S7 and, in spirit, in M4.]}
\item Reproduce Theorem \ref{thm:farkas}, writing the ``not both'' argument separately and in
one line.
\item Prove Corollary \ref{cor:farkasineq} and Theorem \ref{thm:gordan}, and for each identify
the cone to which Theorem \ref{thm:farkas} was applied.
\item Use Theorem \ref{thm:farkas} to complete the Remark after Proposition \ref{prop:polarex}:
show $\{x:Ax\le0\}^\circ=\cone(A^\top)$, with no closure.
\item \textbf{The central exercise of the module.} Write the proof of Theorem \ref{thm:farkas}
on one page so that the appeal to Theorem \ref{thm:sep} is a single displayed step and the
appeal to Lemma \ref{lem:closedcone} is a single marked line. \textbf{Keep this page. M4 repeats
it with the tangent cone and obtains KKT; M11 repeats it with the constraint cone of a linear
program and obtains LP duality. You will be asked in M11 to lay the three proofs side by side,
and they must look like one proof.}
\end{enumerate}

\mpart{Session gate.} Checks 2, 3, 4 and 5 pass; E6.1, E6.2 and E6.5 are written out.

\clearpage

\section{Session S7. Consolidation, capstone and checkpoint}

\mpart{The question.} Three hypotheses have been carried through six sessions: convexity,
closedness, and finite generation. Each was used, and each was claimed to be load-bearing. Is
that true? What exactly goes wrong when each is dropped --- and, since real computation is
inexact, what goes wrong when a hypothesis holds but only barely?

\mpart{The picture.} Two failures, one qualitative and one quantitative.

Qualitatively: there are convex cones that are not closed, and on such a cone the bipolar theorem
fails, strong separation fails, and Farkas' lemma would be false. Lemma \ref{lem:closedcone}
says no \emph{finitely generated} cone can behave this way, which is why Theorem
\ref{thm:farkas} is an alternative. So the qualitative failure is not reachable from a matrix
$A$ --- it needs infinitely many generators, which is what semi-infinite and general convex
programming have, and it is where duality gaps live.

Quantitatively: what \emph{is} reachable from a matrix is near-failure. Make the generators of
$\cone(A)$ nearly dependent and put $b$ just outside. Farkas' lemma remains exactly true and the
certificate remains exactly valid --- but its margin $\dist_{\cone(A)}(b)^2$ shrinks towards
zero, and below the solver's tolerance the two alternatives become indistinguishable
numerically. The theorem is intact and the decision is impossible. \emph{This is the same
phenomenon M4 calls a constraint qualification that holds only barely and M13 calls a badly
centred starting point}, and meeting it once, here, at a moment when nothing depends on it, is
the purpose of the capstone.

\subsection*{The capstone: a convex cone that is not closed}

Let
\[
K=\{(u,v)\in\R^2:u>0\}\cup\{(0,0)\} .
\]
$K$ is a cone: $\lambda(u,v)$ has positive first coordinate for $\lambda>0$ and $u>0$, and
$\lambda=0$ gives the origin. $K$ is convex: a convex combination of two points with $u>0$ has
$u>0$, and combining either with the origin keeps $u\ge0$ with $u>0$ unless all the weight is at
the origin. $K$ is not closed: $(1/k,0)\in K$ and $(0,0)\in K$, but $(0,1)\notin K$ while
$(1/k,1)\in K$ converges to it.

Compute the polar. If $(p,q)\in K^\circ$ then $pu+qv\le0$ for every $u>0$ and every $v\in\R$;
letting $v\to\pm\infty$ forces $q=0$, and letting $u\to0^+$ then $u\to\infty$ forces $p\le0$.
Conversely $(p,0)$ with $p\le0$ is in $K^\circ$. So
\[
K^\circ=\{(p,0):p\le0\},\qquad\text{and hence}\qquad
K^{\circ\circ}=\{(u,v):pu\le0\ \forall p\le0\}=\{(u,v):u\ge0\}=\cl K .
\]
So $K^{\circ\circ}=\cl K\supsetneq K$: Theorem \ref{thm:bipolar} holds, in the form it was
proved --- $S^{\circ\circ}=\cl\conv\cone(S)$ --- and the identity $K^{\circ\circ}=K$ fails
exactly because $K$ is not closed. Moreover the point $b=(0,1)\in\cl K\setminus K$ cannot be
strongly separated from $K$: every $a\in K^\circ$ has $\ip{a}{b}=0$, so no hyperplane puts $b$
strictly on one side. \emph{Farkas' lemma would be false for this cone}, and by Lemma
\ref{lem:closedcone} no matrix can produce it.

\subsection*{The hypothesis table}

Reproduce and complete this table from memory as part of the checkpoint.

\begin{center}
\small
\begin{tabular}{@{}L{3.0cm} L{5.0cm} L{6.6cm}@{}}
\toprule
\textbf{Hypothesis} & \textbf{What it buys} & \textbf{What breaks without it} \\
\midrule
Convexity of $C$
  & Uniqueness in Thm.~\ref{thm:proj}(1); the characterization in Thm.~\ref{thm:proj}(2)
  & The nearest point need not be unique; $\dist_C$ is no longer convex;
    Cor.~\ref{cor:intersection} fails, so there is no dual description \\
Closedness of $C$
  & Existence in Thm.~\ref{thm:proj}(1); strong separation, Thm.~\ref{thm:sep}
  & The infimum is approached only at infinity (Ex.~\ref{ex:notclosed});
    $K^{\circ\circ}\ne K$ (S7 capstone); duality gaps in M4 \\
Compactness of one set
  & Closedness of $C-D$, hence Thm.~\ref{thm:septwo}
  & Disjoint closed convex sets with no strong separation
    (Ex.~\ref{ex:nosep}) \\
Finite generation
  & Lemma \ref{lem:closedcone}, hence Thm.~\ref{thm:farkas} as an \emph{alternative}
  & Non-closed cones exist; the certificate may not; this is what a constraint
    qualification restores in M4 \\
Finite dimension
  & Thm.~\ref{thm:cara}; the compactness of the unit sphere in Thm.~\ref{thm:support}
  & Carath\'eodory is false as stated; supporting hyperplanes need a different argument \\
\bottomrule
\end{tabular}
\end{center}

\subsection*{Proof of concept: the suite, the non-closed cone, and near-degeneracy}

Assemble \texttt{s1.py} to \texttt{s6.py} into one runner that imports \texttt{m1lib.py},
executes every validation check of the module with the single module seed, and prints a table of
check, value, tolerance and verdict.

\emph{Validation.}
\begin{enumerate}
\item Every check from S1 to S6 passes when run from the single runner with the module seed.
\emph{If a check passes only in its own file, something was hard-coded; find it.}
\item For the capstone cone $K$ and $b=(0,1)$: sample unit vectors $a$, keep those with
$\sup_{c\in K}\ip{a}{c}\le0$ (numerically, over a large sample of $c\in K$), and report
$\max\ip{a}{b}$ over them. Confirm it is $0$ to the sampling resolution, so no strong separation
exists; and confirm $b\in K^{\circ\circ}$ using S5's code. \textbf{Report both numbers together:
this is the closedness hypothesis of Theorems \ref{thm:sep} and \ref{thm:bipolar} failing.}
\item Build the near-degenerate family
$A_\varepsilon=\begin{bmatrix}1&1\\0&\varepsilon\end{bmatrix}$ and
$b_\varepsilon=(1,-\varepsilon)$. The two generators $(1,0)$ and $(1,\varepsilon)$ span the thin
wedge $\cone(A_\varepsilon)=\{(u,v):v\ge0,\ u\ge v/\varepsilon\}$, and $b_\varepsilon$ sits just
below it, so system (I) is infeasible for \emph{every} $\varepsilon>0$. Here
$\proj_{\cone(A_\varepsilon)}b_\varepsilon=(1,0)$ and the Farkas margin is
$\norm{b_\varepsilon-\proj b_\varepsilon}^2=\varepsilon^2$: the certificate
$y=(0,\varepsilon)$ is exactly valid at every $\varepsilon$, and its margin decays
\emph{quadratically}. For $\varepsilon=10^{-1},\dots,10^{-6}$, tabulate the computed margin
against $\varepsilon^2$ and confirm they agree to solver tolerance until the margin falls below
it. Report the smallest $\varepsilon$ at which your feasibility solve and your certificate
solve disagree about the instance --- with a default \texttt{linprog} tolerance this happens
well before $\varepsilon=10^{-6}$. \textbf{Keep this table: the theorem is exactly true at every
$\varepsilon$ and the numerical decision becomes impossible anyway. M13 asks why an interior
point method needs a well-centred start, and this is the answer.}
\end{enumerate}

\emph{Deliverable.} \texttt{run\_m1.py}, its printed verdict table with one row per check for
the whole module, and the two tables of checks 2 and 3.

\subsection*{Checkpoint}

The module is complete when the following can be produced from scratch, on paper, with this
document closed:

\begin{enumerate}
\item Theorem \ref{thm:proj}, all three parts, with proof, and with the line using closedness
and the line using convexity marked.
\item Theorem \ref{thm:sep}, with proof, the separating vector written explicitly as
$x-\proj_Cx$, and the margin identified as $\dist_C(x)^2$.
\item Lemma \ref{lem:closedcone} and Theorem \ref{thm:farkas}, with proofs, the single line
where closedness is used marked, and the one-line ``not both'' argument written separately.
\item Theorem \ref{thm:moreau}, with proof, and the sentence identifying the orthogonal
decomposition of a subspace as its special case.
\item The hypothesis table above, reproduced and with at least one entry of your own added.
\end{enumerate}

\mpart{Keep item 3.} M4 asks the identical question with the tangent cone of the active
constraints in place of $\cone(A)$, and the answer is the Karush--Kuhn--Tucker theorem; M11 asks
it with the constraint cone of a linear program, and the answer is LP duality. The method is
Theorem \ref{thm:farkas} exactly as used here; only the cone changes. Without this proof, KKT
will look like a list of conditions someone wrote down, and LP duality like a separate theory.

\mpart{Keep item 1.} M8 uses the fixed-point equation $x=\proj_C(x-t\nabla f(x))$ as its
stationarity condition, and it is Theorem \ref{thm:proj}(2) restated; M9 uses Theorem
\ref{thm:proj}(3), the firm version, to prove Fej\'er monotonicity of proximal iterates. Both
will look like new material if the projection theorem is not in hand.

\mpart{Session gate.} The runner is green and the checkpoint has been written out on paper. M1
is finished; M2 may be started.

\clearpage

\section{Common errors}

\begin{itemize}
\item \textbf{Reading ``closed'' as boilerplate.} It is the hypothesis that fails, and S7 is
built around the failure. Strong separation, the bipolar theorem and Farkas' lemma are all false
without it, and every duality gap in the plan is this failure in disguise.

\item \textbf{Confusing the interior with the relative interior.} $\Simplex_n$ has empty
interior in $\R^n$ (Example \ref{ex:simplexri}). Any argument that produces an interior point of
a set defined partly by equalities has a bug; the correct statement uses $\ri$.

\item \textbf{Proving the projection formula by solving the optimization problem.} Once Theorem
\ref{thm:proj}(2) is available, every closed form in Proposition \ref{prop:easyproj} and
\ref{prop:simplex} is verified by \emph{checking an inequality on a guess}. If a derivation in
your notes optimizes, it is longer than it needs to be and will not generalize.

\item \textbf{Expecting $\proj_C$ to be linear.} It is nonexpansive, not linear, unless $C$ is
affine. The consequence in M8 is that a projected gradient step is not a gradient step followed
by a separable correction. S3 check 5 exists to make this a number.

\item \textbf{Clipping and renormalizing instead of projecting onto the simplex.} The worked
example of S3 gives a discrepancy of $0.248$ on a four-dimensional instance. This is the single
most common implementation error in projected policy-gradient code.

\item \textbf{Saying ``separating hyperplane'' without saying which separation.} Proper, strong
and supporting are three theorems with three hypotheses. Which one is available determines
whether the duality gap of M4 is zero. E4.5 exists for this.

\item \textbf{Checking that a separation occurred rather than checking the margin.} Theorem
\ref{thm:sep} predicts the gap exactly: $\dist_C(x)^2$. An implementation that separates with
the wrong margin has a bug that the weaker check will never find. S4 check 1 and S6 check 3 are
the two places this matters.

\item \textbf{Mixing the polar and the dual cone.} $C^*=-C^\circ$. Both conventions are in
print; this document uses the polar throughout. A sign error here propagates silently into every
KKT condition in M4. Assert the convention in a docstring in \texttt{m1lib.py}.

\item \textbf{Believing both Farkas alternatives could fail.} They cannot, and the reason is
Lemma \ref{lem:closedcone} and nothing else. If an implementation ever reports neither, the
tolerance in the certificate check is looser than the tolerance in the feasibility solve --- a
numerical artifact of exactly the kind S7 check 3 is about, not a counterexample.

\item \textbf{Treating S1 as the soft session.} Corollary \ref{cor:conecara} is the proof of
Lemma \ref{lem:closedcone}, which is the hypothesis of Theorem \ref{thm:farkas}, which is the
theorem the rest of the plan quotes. The module's logical chain runs through its shortest
session.
\end{itemize}

\section{What M1 hands forward}

\begin{table}[htbp]
\centering
\small
\caption{Results established in M1, the session that establishes each, and where each is used
again.}
\label{tab:forward}
\begin{tabular}{@{}L{4.7cm} l L{1.8cm} L{6.0cm}@{}}
\toprule
\textbf{Result} & \textbf{From} & \textbf{Used in} & \textbf{As} \\
\midrule
Projection theorem (Thm.~\ref{thm:proj})
  & S3 & M8
  & Stationarity as $x=\proj_C(x-t\nabla f(x))$, and the gradient mapping \\
Projection theorem, part 3
  & S3 & M9
  & Firm nonexpansiveness, hence Fej\'er monotonicity of proximal iterates \\
Projection theorem, affine case
  & S3 & MMC201 M2
  & The normal equations; least squares as an orthogonal projection \\
$\tfrac12\dist_C^2$ and $\proj_C$ (Prop.~\ref{prop:dist})
  & S3 & M3
  & The Moreau envelope and proximal operator of the indicator of $C$ \\
Simplex projection (Prop.~\ref{prop:simplex})
  & S3 & M8, M9, M15
  & The projected policy-gradient step, at $O(n\log n)$ per state \\
Strong separation (Thm.~\ref{thm:sep})
  & S4 & M4
  & Strong duality under Slater: separate $(0,p^\star)$ from the perturbation epigraph \\
Failure of strong separation (Ex.~\ref{ex:nosep})
  & S4 & M4
  & What a nonzero duality gap is \\
Supporting hyperplane (Thm.~\ref{thm:support})
  & S4 & M11, M14
  & Faces and extreme points of a polyhedron; cutting planes \\
Proper separation (Thm.~\ref{thm:proper})
  & S4 & M4
  & Why Slater's condition is stated with $\ri$ \\
Bipolar theorem (Thm.~\ref{thm:bipolar})
  & S5 & M4, M11
  & Dual descriptions of cones; the dual of a polyhedral constraint set \\
Moreau decomposition (Thm.~\ref{thm:moreau})
  & S5 & M3
  & The conic case of $x=\prox_f(x)+\prox_{f^*}(x)$ \\
Normal cone (Def.~\ref{def:normal}, Prop.~\ref{prop:normalsimplex})
  & S5 & M4, M8
  & $-\nabla f(x^\star)\in\Nc_C(x^\star)$ as constrained stationarity \\
Closedness of $\cone(A)$ (Lem.~\ref{lem:closedcone})
  & S6 & M11
  & Why linear programming needs no constraint qualification \\
Farkas' lemma (Thm.~\ref{thm:farkas})
  & S6 & M4
  & KKT, with the tangent cone in place of $\cone(A)$ \\
Farkas' lemma
  & S6 & M11, M12
  & LP duality; the Phase I infeasibility certificate \\
Farkas' lemma
  & S6 & M14
  & Benders feasibility cuts as extreme rays of the dual subproblem \\
Fourier--Motzkin (E2.4)
  & S2 & M14
  & Projecting out the complicating variables in Benders decomposition \\
Carath\'eodory (Thm.~\ref{thm:cara})
  & S1 & M8, M11
  & Sparsity of Frank--Wolfe iterates; mixtures of deterministic policies \\
\bottomrule
\end{tabular}
\end{table}

\section{References}

Nothing in this list is required reading: every proof above is written out above. Each entry
says what it adds.

\begin{enumerate}[label={[\arabic*]}]
\item Hiriart-Urruty, J.-B. and Lemar\'echal, C. (2001). \emph{Fundamentals of Convex
Analysis}. Springer. Part A. The closest published treatment to this module's sequence, with
more examples and the infinite-dimensional caveats.
\item Boyd, S. and Vandenberghe, L. (2004). \emph{Convex Optimization}. Cambridge University
Press. Ch.~2 and \S 8.1: a large catalogue of convex sets and closed-form projections. Freely
available.
\item Crouzeix, J.-P. and Oca\~na, E. (2018). \emph{An\'alisis convexo}. UNI. The MMC007
syllabus text; use it to fix the Spanish vocabulary of the lectures.
\item Rockafellar, R. T. (1970). \emph{Convex Analysis}. Princeton University Press. \S 6, 11,
14, 17, 19, 22. The reference statements later modules cite, and the source of the two proofs
omitted here (\S 6 for Prop.~\ref{prop:ricl}(2), Thm.~11.3 for Thm.~\ref{thm:proper}).
\item Bertsekas, D. P. (1999). \emph{Nonlinear Programming}, 2nd ed. Athena Scientific.
App.~B: recession cones and closedness criteria in far more detail than S2 needs.
\item Bauschke, H. and Combettes, P. (2017). \emph{Convex Analysis and Monotone Operator Theory
in Hilbert Spaces}, 2nd ed. Springer. The operator reading of $\proj_C$ as a firmly
nonexpansive map, which is the language M9 will use.
\item Schrijver, A. (1986). \emph{Theory of Linear and Integer Programming}. Wiley. \S 7.3:
Farkas' variants collected, proved polyhedrally rather than by separation --- worth seeing once
as a contrast with S6.
\item Bertsimas, D. and Tsitsiklis, J. (1997). \emph{Introduction to Linear Optimization}.
Athena Scientific. \S 4.6: Farkas read as an infeasibility certificate, the form M11 and M12
use.
\item Beck, A. (2017). \emph{First-Order Methods in Optimization}. SIAM. Ch.~6: a larger
projection table, and the projection as a special proximal operator, which is M3.
\item Farkas, J. (1902). Theorie der einfachen Ungleichungen. \emph{J.\ Reine Angew.\ Math.}
124, 1--27. Theorem \ref{thm:farkas} in the original.
\item Gordan, P. (1873). \"Uber die Aufl\"osung linearer Gleichungen mit reellen Coefficienten.
\emph{Math.\ Ann.} 6, 23--28. Theorem \ref{thm:gordan} in the original.
\item Stiemke, E. (1915). \"Uber positive L\"osungen homogener linearer Gleichungen.
\emph{Math.\ Ann.} 76, 340--342; and Motzkin, T. (1936). \emph{Beitr\"age zur Theorie der
linearen Ungleichungen}. Dissertation, Basel. The two alternatives named but not proved in S6.
\item Moreau, J.-J. (1962). D\'ecomposition orthogonale d'un espace hilbertien selon deux
c\^ones mutuellement polaires. \emph{C.\ R.\ Acad.\ Sci.\ Paris} 255, 238--240. Theorem
\ref{thm:moreau} in the original; two pages.
\item Duchi, J., Shalev-Shwartz, S., Singer, Y. and Chandra, T. (2008). Efficient projections
onto the $\ell_1$-ball for learning in high dimensions. \emph{ICML}; and Condat, L. (2016). Fast
projection onto the simplex and the $\ell_1$ ball. \emph{Mathematical Programming} 158,
575--585. The $O(n\log n)$ algorithm of Proposition \ref{prop:simplex} and its refinements.
\item Zhao, S. (2025). \emph{Mathematical Foundations of Reinforcement Learning}. Springer. The
expository method this document follows --- question, picture, statement, proof --- and the
natural destination of the reinforcement-learning remarks in each session.
\end{enumerate}

\end{document}
